\begin{corollary}[E.2.4 (Proved): arbitrary connected deformation chains]
The ordered coprime pair $(3k,4l)$ is constant in every
proper smooth holomorphic family over a connected complex
analytic base which contains a constructed sphere member.
Two constructed members are deformation equivalent through
such families, allowing arbitrary intermediate members and
finite chains, if and only if their pairs agree.
Moreover different pairs cannot occur in one continuous
family in the $C^\infty$ space of integrable complex
structures on $S^6$.
\end{corollary}
\begin{proof}
The detailed proof is E.2.5--E.2.24 below.  The intrinsic
pencil persists for either canonical slope, the limiting
period stays in the admissible interior, and the two finite
root charts and the actual analytic cusp are recovered.
E.2.21 then identifies every full punctured-disc smooth
completion with the prescribed member.  The additional
step E.2.22 propagates from any nonempty known open subset
of a disc: finitely many successive discrete analytic
exceptional sets are filled in reverse order.  E.2.23 uses
a finite local projection and normalization of a curve
through two specified points to pass to arbitrary analytic
bases and deformation chains.
The classical $C^\infty$ Kuranishi local theorem then
gives the continuous-family conclusion in E.2.24.

Local miniversality alone would prove only openness.
Nor does a boundary point of an arbitrary open locus
necessarily admit a holomorphic disc whose punctured
image lies in that locus.  The generic-disc propagation
and finite-projection steps are the substantive additions
which establish the required closedness.
\end{proof}

% Owner: boundary_d.  Separate from the E-owned master source.
% These auxiliary assertions concern one-sided smooth specializations.
% The full deformation-chain conclusion is proved in E.2.22--E.2.24.
\subsubsection*{Smooth specializations of the canonical pencil}

\begin{proposition}[E.2.5 (Proved): a fractional canonical pencil at a smooth limit]
Let $\pi:\mathcal X\to\Delta$ be a proper holomorphic submersion
whose fibres are diffeomorphic to $S^6$.  Suppose that for every
$t\ne0$ the fibre is one of the members of E.2.1, with the same
coprime multiplicities $(3k,4l)$.  Set
\[
 A=12kl,\qquad D=kl-3k-4l\ne0.
\]
There is a unique holomorphic line bundle $\mathcal L$ on
$\mathcal X$, up to isomorphism, satisfying
\[
 \mathcal L^{D}=K_{\mathcal X/\Delta}^{A}.
\]
Its restriction to $X_t$ for $t\ne0$ is the coarse pencil
$f_t^*\mathcal O_{\mathbb P^1}(1)$.  The coherent sheaf
$\pi_*\mathcal L$ is locally free of rank two, and its central
fibre injects into $H^0(X_0,\mathcal L_0)$.

The central manifold has algebraic dimension one and a
holomorphic algebraic reduction $g_0:X_0\to\mathbb P^1$ with
connected fibres.  If $s_0,s_1$ are the two central sections
obtained from a local frame of $\pi_*\mathcal L$, there are a
nonconstant rational map $R:\mathbb P^1\to\mathbb P^1$ of
degree $e\ge1$ and an effective Cartier divisor $E_0$ such that
\[
 [s_0:s_1]=R\circ g_0,\qquad
 \mathcal L_0\simeq g_0^*\mathcal O(e)\otimes\mathcal O(E_0).
\]
Here $[s_0:s_1]$ denotes the holomorphic extension of the ratio;
the displayed divisor is its common fixed divisor.  The
statement does not assert that $E_0=0$.
\end{proposition}
\begin{proof}
By Ehresmann's theorem the smooth bundle over the contractible
disc is differentiably trivial.  Consequently
$H^1(\mathcal X,\mathbb Z)=H^2(\mathcal X,\mathbb Z)=0$.
The exponential sequence identifies
\[
 \operatorname{Pic}(\mathcal X)=H^1(\mathcal X,\mathcal O_{\mathcal X})
\]
as groups.  This is a complex vector space, so multiplication by
the nonzero integer $D$ is an automorphism.  This proves the
existence and uniqueness of the \emph{relative} root, including
when $D<0$.  The same argument on each fibre makes its Picard
group uniquely divisible.  E.2.1 gives
$K_{X_t}^{A}=(f_t^*\mathcal O(1))^{D}$, so uniqueness identifies
$\mathcal L_t$ with this actual pencil, not just with a numerical
class.

Proper direct-image coherence gives $\mathcal V=\pi_*\mathcal L$.
It is torsion-free: a local section killed by a nonzero function
of $t$ vanishes on the inverse image of an open set, hence
vanishes as a section of the invertible sheaf on the smooth
total space.  A coherent torsion-free sheaf on a smooth curve
is locally free.  Its rank is two because the noncentral
algebraic reductions have connected fibres.  Applying $\pi_*$
to
\[
 0\longrightarrow\mathcal L\xrightarrow{\,t\,}\mathcal L
 \longrightarrow i_*\mathcal L_0\longrightarrow0
\]
shows directly that $\mathcal V/t\mathcal V$ injects into
$H^0(X_0,\mathcal L_0)$.  Thus the two central sections are
linearly independent even if central cohomology jumps; no
surjectivity of base change has been assumed.

We use precisely the non-holomorphic-meromorphic-map theorem
of \citet[Theorem~2.1]{CDP2020D}, also written out in
\citet[Theorem~3.1 and Corollary~3.2]{LRS2018D}.
For a compact complex threefold with $b_2=0$, a nonconstant
meromorphic map to a curve which fails to be holomorphic
would force $c_3=0$.  Since $c_3(X_0)[X_0]=2$, every
meromorphic function on $X_0$ extends as a holomorphic map
to $\mathbb P^1$.  This argument uses neither the different
holomorphic-fibration assertion of CDP, Theorem~2.2, nor its
Corollary~2.3.  Those assertions cannot be applied here.

In particular the ratio of $s_0,s_1$ gives a nonconstant
holomorphic map $h_0:X_0\to\mathbb P^1$.  Also $a(X_0)\le1$:
otherwise two algebraically independent meromorphic functions
would give a holomorphic map $X_0\to\mathbb P^1\times\mathbb P^1$
with two-dimensional image.  At a point of differential rank
two, after choosing finite affine target coordinates $u,v$
centred at its image, the global meromorphic function $u/v$
would have the local indeterminacy of the quotient of two
independent coordinates.  It cannot extend holomorphically
there.  This contradicts the preceding theorem.  Hence
$a(X_0)=1$.

Take the Stein factorization $h_0=R\circ g_0$.  Its intermediate
space is a smooth compact curve.  Connectedness of the fibres
injects its $H^1(\ ,\mathbb Z)$ into $H^1(X_0,\mathbb Z)=0$
by the degree-one Leray sequence, so it is $\mathbb P^1$.
Every meromorphic function on $X_0$ factors rationally through
$g_0$: the image of its holomorphic map together with $g_0$
is an analytic curve in $\mathbb P^1\times\mathbb P^1$;
normalizing this curve gives a finite factor of $g_0$, whose
degree is one because $g_0$ has connected fibres.  Thus $g_0$
is the algebraic reduction, and $R$ has some degree $e\ge1$.

Finally, on a target chart where the first homogeneous
coordinate of $h_0$ is nonzero, $s_1=h_0s_0$ in a local
trivialization, and the common ideal is generated by $s_0$.
The analogous assertion holds on the other chart.  These local
generators define a nonzero section of
$\mathcal L_0\otimes h_0^*\mathcal O(-1)$ and an effective
Cartier zero divisor $E_0$.  This gives the stated factorization
of $\mathcal L_0$ with no residual codimension-two base locus.
\end{proof}

\begin{proposition}[E.2.6 (Proved): positive canonical slope preserves the fibration]
Under the hypotheses of E.2.5, suppose $D>0$.  Then $E_0=0$
and $e=1$.  After shrinking the disc and choosing a frame of
$\pi_*\mathcal L$, its evaluation defines a proper flat
holomorphic map
\[
 g:\mathcal X\longrightarrow\mathbb P^1\times\Delta
\]
which is the intrinsic algebraic reduction on every fibre.
Its smooth general central fibre is a complex two-torus.

In particular, two families of the type in E.2.5 with positive
canonical slope cannot have the same smooth central manifold
unless their multiplicity pairs agree.  A positive-slope and
a negative-slope family cannot have a common smooth central
manifold either.  This proposition concerns a common smooth
specialization.  The additional propagation and closure
arguments for arbitrary connected chains are E.2.22--E.2.24.
\end{proposition}
\begin{proof}
Choose a very general smooth fibre $F$ of $g_0$, avoiding the
images of all vertical components of $E_0$ and the jumping
sets for the pluricanonical direct images below.  Its normal
bundle in $X_0$ is trivial.  Adjunction gives
$K_F=K_{X_0}|_F$, and $c_1(K_F)=0$ integrally because
$H^2(X_0,\mathbb Z)=0$.  In particular $F$ is minimal: a
$(-1)$-curve would have canonical intersection $-1$.
Restriction of the root identity in E.2.5 gives
\[
 K_F^{A}\simeq\mathcal O_F(D E_0|_F).
\]
The right hand side is effective since $D>0$, so
$\kappa(F)\ge0$.

We next show $\kappa(F)\le0$, without assuming that the fibres
of an algebraic reduction have algebraic dimension zero.
If $\kappa(F)>0$, some $r>0$ has $h^0(F,K_F^r)\ge2$.
Our very-general choice ensures that this is the generic rank
of the coherent sheaf $g_{0*}K_{X_0}^r$.  A sufficiently
positive twist of this sheaf on $\mathbb P^1$ is globally
generated.  Two of its global sections can therefore be
chosen to be linearly independent on a general fibre.  By
projection formula they are sections of
$K_{X_0}^r\otimes g_0^*\mathcal O(N)$; their ratio is a
meromorphic function on $X_0$ which is nonconstant on $F$.
This contradicts the algebraic-reduction assertion of E.2.5.
The choice of a very general fibre is legitimate: for each
$r$ the exceptional set for the generic cohomology rank is
a proper analytic subset of the base curve, and there are
only countably many $r$.

Thus $\kappa(F)=0$.  The classification of minimal compact
complex surfaces makes $K_F$ torsion; the relevant consequence
is also explicitly recorded in the proof of
\citet[Proposition~3.1.2]{AbashevaDeev2024AB}.
An effective divisor whose line bundle is torsion is zero:
raising its canonical section to a torsion-killing power
gives a holomorphic function on a compact connected manifold,
and that function is a nonzero constant.  Therefore
$E_0|_F=0$.  Every horizontal component of $E_0$ would meet
such a general fibre in a nonempty effective divisor.
Consequently $E_0$ has only vertical components, supported
over finitely many points of $\mathbb P^1$.

We now use connectedness of actual divisor limits to remove
both this vertical fixed part and the possible degree $e$.
Choose a linear combination $s=as_0+bs_1$ for which its
moving central divisor is the union of the $e$ distinct
smooth fibres over $R^{-1}([-b:a])$, all disjoint from the
support of $E_0$.  Also choose it to be a regular fibre on
one specified noncentral member.  Both requirements omit
only finitely many choices in the projective line of
combinations.  Extend it by the corresponding combination
of the frame of $\pi_*\mathcal L$.

Its total zero divisor $\mathcal D\subset\mathcal X$ is a
relative effective Cartier divisor: the section is not
identically zero on any fibre, by independence of the two
sections, including at zero.  In local coordinates, $t$ is
a non-zero-divisor modulo its defining equation.  Hence
$\mathcal D\to\Delta$ is proper and flat.  Its general
noncentral fibre is a smooth connected torus.  There are no
vertical irreducible components of $\mathcal D$, since such
a component would have to be the whole irreducible
threefold $X_t$.  It follows that $\mathcal D$ is reduced
and irreducible: its unique generic component has multiplicity
one, and a Cartier divisor in a smooth manifold has no
embedded components.  Normalize $\mathcal D$ and take the
Stein factorization over $\Delta$.  The finite factor has
degree one over the punctured generic locus; normality of
the disc makes it the disc itself.  Its central fibre, and
therefore the image $\mathcal D_0$, is connected.

But the support of
\[
 \mathcal D_0=E_0+g_0^*R^*([-b:a])
\]
is a disjoint union of the $e$ chosen smooth connected fibres
and, if nonempty, the support of $E_0$ over other base points.
Connectedness forces $E_0=0$ and $e=1$.

The central two sections now generate $\mathcal L_0$ at
every point.  The noncentral two sections already generate
their pencils, so the evaluation
$\pi^*\pi_*\mathcal L\to\mathcal L$ is surjective after
shrinking.  It defines the asserted map $g$.  It is proper
because $\pi$ is proper.  On each $X_t$ the nonconstant
map to the smooth curve $\mathbb P^1$ is flat: its local
structure sheaf is torsion-free over the local discrete
valuation ring of the curve.  The total source is flat
over $t$, so the local criterion for flatness applied to
the parameter $t-t_0$ gives flatness over $(z,t)$.
Equivalently, the local fibre equations form a regular
sequence in the smooth source.  At a regular
point of the central base curve, properness makes $g$ a
proper submersion on a neighbourhood of the corresponding
compact fibre.  Ehresmann's theorem identifies that central
surface differentiably with nearby complex two-tori.
The classification of compact complex surfaces, or its
special case for the differentiable four-torus, then makes
the central surface a complex two-torus as well.

Suppose now two positive-slope families have the same
central manifold $X_0$.  Its algebraic reduction is intrinsic
up to $\operatorname{PGL}_2$, so the preceding result makes
both central pencil lines the same non-torsion line
$L=g_0^*\mathcal O(1)$.  With indices $i=1,2$ the identities
$K_{X_0}^{A_i}=L^{D_i}$ give
$L^{D_1A_2-D_2A_1}=\mathcal O$.  The line $L$ is non-torsion
because a positive power has a nonzero section with zeros.
Thus $D_1/A_1=D_2/A_2$, and the exact recovery of the ordered
pair from this slope in E.2.3 gives the assertion.

Finally, a common smooth limit from opposite signs would
have a nonzero section of a positive power of $K_{X_0}$
and of a negative power, by upper semicontinuity (or by
the same injective direct-image argument).  Raising them
to a common power, their nonzero product is a constant;
both are nowhere zero.  Hence $K_{X_0}$ is torsion.
The exponential sequence and $H^1(X_0,\mathbb Z)=
H^2(X_0,\mathbb Z)=0$ make $\operatorname{Pic}(X_0)$
torsion-free, so $K_{X_0}$ would be trivial.  Its nonzero
holomorphic three-form $\Omega$ is closed and exact because
$H^3(X_0,\mathbb C)=0$.  Stokes' theorem would give
$\int_{X_0}\Omega\wedge\overline\Omega=0$, whereas its
appropriate constant multiple is a positive volume form.
This contradiction excludes the common limit.
\end{proof}

\begin{remark}[Scope after E.2.6]
E.2.5 proves that the central meromorphic reduction is
holomorphic, and E.2.6 removes its fixed divisor and degree
loss for $D>0$.  The positive-slope proof alone does not
exclude the effective antipluricanonical divisor allowed
by the restriction identity when $D<0$.  E.2.8--E.2.10
remove that possibility by the horizontal-divisor vanishing
argument.  E.2.13 and E.2.16--E.2.21 then recover the finite
charts, cusp, and actual smooth completion.  The separate
generic-disc propagation and closure arguments in
E.2.22--E.2.24 prove E.2.4 for arbitrary analytic bases,
deformation chains, and continuous complex structures.
The present proposition retains its stated one-sided,
positive-slope scope.
\end{remark}

% Owner: classification_e. Root and boundary_d full-source referee PASS.
% This auxiliary proposition controls the period; E.2.20--E.2.24 complete the classification.
\begin{proposition}[E.2.7 (Proved): an extended fibration cannot reach the zero period]
Under the hypotheses of E.2.5, suppose in addition that the
intrinsic algebraic reductions extend to a holomorphic map
\[
 g:\mathcal X\longrightarrow\mathbb P^1\times\Delta
\]
whose restriction to $X_0$ is nonconstant.  This hypothesis
holds in E.2.6.  Normalize the punctured family to fixed
discrete logarithmic data with multiplicities $(3k,4l)$, and
write its isomorphism parameter as
\[
 z=\exp\!\left(-\frac{\pi i c}{kl}\right).
\]
The resulting holomorphic function $z:\Delta^*\to\mathbb C^*$
extends across zero and satisfies $z(0)\ne0$.  Its central
value lies strictly inside the admissible period disc.

This statement concerns the period function and the extended
relative fibration.  Identification of $X_0$ with the member
at $z(0)$ is proved separately in E.2.20--E.2.21; propagation
to arbitrary deformation bases and continuous families is
proved in E.2.22--E.2.24.
\end{proposition}
\begin{proof}
Local miniversality in E.2.1 supplies holomorphic local
classifying functions for the family on $\Delta^*$.
E.1.2--E.1.3 normalize the discrete data by holomorphic
family comparisons.  For fixed data the comparison shifts
are precisely $c\mapsto c+2kl\mathbb Z$, so the displayed
exponential is a single-valued holomorphic function on
$\Delta^*$.  An automorphism with the same marked parameter
has identity linear action on the homology of an ordinary
torus: in the necessity argument of E.1.2 its parameters
are $\alpha=1$ and $h=0$.  Consequently the homology local
system of the punctured family is also the one obtained
from these local classifying maps; possible translation
automorphisms of the comparisons add no homology monodromy.

The admissible parameter domain is a punctured disc
$0<|z|<R_{k,l}$, with finite positive radius.  The removable
singularity theorem therefore extends $z$ holomorphically
to $\Delta$.  It remains to exclude a zero at the centre.

We first control the moving three critical values without
assuming that they remain distinct at zero.  Since
$\pi:\mathcal X\to\Delta$ is a submersion, the critical
locus of $g$ is the analytic locus where its differential
has rank less than two.  Properness makes its image
\(\mathcal C\subset\mathbb P^1\times\Delta\) a closed
analytic set.  Its intersection with the central target
curve is the finite critical-value set of the nonconstant
map $g_0:X_0\to\mathbb P^1$.  In particular
$\mathcal C$ has no component equal to that whole curve.
After shrinking $\Delta$, its reduced one-dimensional
part is finite over $\Delta$ and over $\Delta^*$ consists
of the three critical-value sections.

Those sections are individually marked by the linear
peripheral monodromy orders $3$, $4$, and infinity.
The normalization of the closure of each section is
therefore finite of degree one over the disc.  It is
isomorphic to the disc, so the three values
$c_1(s),c_2(s),c_0(s)$ extend as holomorphic maps to
$\mathbb P^1$, with collisions allowed at $s=0$.
The unique projective transformation $A_s$ sending this
ordered triple to $(0,1,\infty)$ has meromorphic
coefficients at zero, by its explicit cross-ratio formula.
This is the critical-value argument of B.4.4 in a setting
where no vertical discriminant component is present.

Choose a point $v\in\mathbb P^1$ outside all central
critical values.  Properness and closedness of
$\mathcal C$ allow us to shrink the disc so that
$(v,s)\notin\mathcal C$ for every $s$.  Then
\[
 \mathcal T=g^{-1}(\{v\}\times\Delta)\longrightarrow\Delta
\]
is a proper smooth holomorphic family of surfaces.
Its noncentral fibres are ordinary complex two-tori, and
Ehresmann's theorem makes its integral first-homology
monodromy around $s=0$ the identity.

In the original marked base, this same family is obtained
by evaluating the period construction at
\[
 t(s)=A_s(v)\in\mathbb P^1\setminus\{0,1,\infty\}.
\]
The function $t(s)$ is meromorphic at zero.  Its loop on a
sufficiently small positive $s$-circle is therefore either
null-homotopic in the marked base or conjugate to a power
of one of its three peripheral loops.  Indeed, if its
limit is ordinary, its image lies in an ordinary small
disc; if its limit is a marked point, a local coordinate
has a zero or pole of finite order and its image lies
in the corresponding punctured disc.  This observation
also covers a constant $t(s)$.

Let $\Lambda$ have ordered basis
$(\widehat\gamma,\widehat u,\widehat w,\widehat\delta)$,
put $E=\widehat\delta\otimes\gamma$, and write
$S_a=I+aE$.  The two finite peripheral monodromies are
\[
 A_1=\begin{pmatrix}
 1&0&0&0\\6&0&1&0\\-6&-1&-1&0\\-2&1&0&1
 \end{pmatrix},\qquad
 A_2=\begin{pmatrix}
 1&0&0&0\\0&0&-1&0\\-6&1&0&0\\3&0&1&1
 \end{pmatrix}.
\]
They have orders three and four, respectively, even for
the larger multiplicities: the extra finite deck
transformations are translations and do not affect
$H_1$.  In the cusp marking the third monodromy is
\[
 A_0=I+N,\qquad
 N\widehat u=\widehat w,\quad
 N\widehat\gamma=-\widehat\delta,\quad
 N\widehat w=N\widehat\delta=0.
\]
All peripheral matrices preserve $\widehat\delta$ and
$\gamma$, so commute with every $S_a$.  The change to
the cusp marking preserves the same rank-one operator
$E$; conjugating a peripheral loop consequently does
not alter the cancellation test below.

Suppose that $z(s)$ has a zero of order $d>0$ at zero.
Continuation of a logarithm around the positive loop
changes $c$ by $-h$, where
\[
 h=2kl\,d>0.
\]
Since $\Pi_{c+h}=\Pi_cS_h$, a comparison of the identical
complex universal covers from parameter $c-h$ back to
$c$ acts on the marked lattice by $S_{-h}$.
The affine translation part of a closing comparison
acts trivially on $H_1$.  Thus the return map of the
family $\mathcal T$ has homology action, up to a common
conjugation,
\[
 S_{-h},\qquad S_{-h}A_1^n,\qquad
 S_{-h}A_2^n,\qquad\hbox{or}\qquad S_{-h}A_0^n
\]
for an integer $n$.  Reversing the convention for a
peripheral generator only replaces $n$ by $-n$.
The parameter factor and the base factor commute, so
the order in this display has no effect.

None of these matrices is the identity.  In the two
finite cases, taking the third or fourth power would
give $S_{-3h}=I$ or $S_{-4h}=I$, impossible because
$E\ne0$ and $h>0$.  In the cusp case, $N^2=EN=NE=0$
and hence
\[
 S_{-h}A_0^n=I+nN-hE.
\]
Its coefficient of $\widehat w$ on $\widehat u$ is
$n$, so identity would force $n=0$, and then $h=0$.
The ordinary case is $S_{-h}\ne I$.  This contradicts
the identity monodromy of the smooth family over
the full disc.  Therefore $z(0)\ne0$.

Finally, the extended holomorphic function satisfies
$|z(0)|\le R_{k,l}$.  Equality would make it a constant
of maximal modulus, contradicting $|z(s)|<R_{k,l}$
for $s\ne0$.  Thus $0<|z(0)|<R_{k,l}$, as claimed.
The matrix identities and the factor $h=2kld$ are
checked exactly in \path{code/e_push_E_zero_end.py}.
\end{proof}

\begin{remark}[Scope and subsequent completion]
For either canonical slope, E.2.10 supplies the
fibration-persistence hypothesis of this proposition,
so a one-sided smooth specialization has an interior
period value.  Equality of that value with an existing
model does not by itself prove uniqueness of its smooth
completion.  E.2.11--E.2.12 exclude critical values
appearing only at the central parameter; E.2.13--E.2.20
recover the actual finite and cusp charts, and E.2.21
proves the full one-sided identification.  The distinct
problem of propagation through arbitrary deformation
branches is settled by E.2.22--E.2.24.
\end{remark}

% Owner: classification_e.  This is a restriction on the negative-slope
% obstruction, not an exclusion of all Hopf-surface fibrations.
\subsubsection*{A restriction on the negative-slope fixed divisor}

\begin{proposition}[E.2.8 (Proved): horizontal divisors force non-K\"ahler fibre normalizations]
Let $X$ be a compact complex threefold diffeomorphic to $S^6$, and
let $f:X\to\mathbb P^1$ be a surjective holomorphic map with
connected fibres.  Suppose that a smooth general fibre $F$ has
$b_1(F,\mathbb Q)=1$.  Then every reduced fibre of $f$ is
irreducible, and the monodromy on $H^1(F,\mathbb Q)$ is trivial.

If $H\subset X$ is a horizontal prime divisor, the normalization
of every reduced fibre lies outside Fujiki's class $\mathcal C$.
Equivalently, no such normalization admits a K\"ahler resolution.
In particular no fibre has Moishezon normalization.

Without the hypothesis $b_1(F)=1$, every irreducible component
of a \emph{reducible} fibre of any such connected map from a
complex $S^6$ also has normalization outside $\mathcal C$.

Consequently, in E.2.5 with $D<0$, a horizontal component of
$E_0$ is possible only if every reduced fibre of the central
algebraic reduction $g_0$ is irreducible and has normalization
outside $\mathcal C$.  This necessary condition does not exclude
non-K\"ahler normalizations and does not prove persistence of the
relative pencil in the negative-slope case.
\end{proposition}
\begin{proof}
Choose a nonempty finite set $\Sigma\subset\mathbb P^1$ containing
all critical values; a regular value can be added if necessary.
Write $s=|\Sigma|$, put $A=f^{-1}(\Sigma)_{\mathrm{red}}$,
and let $r$ be the number of irreducible components of $A$.
The compact analytic set $A$ is triangulable.  Each of its
irreducible two-dimensional components has its complex
fundamental class, and their intersections have real dimension
at most two.  Hence
\[
 H^4(A,\mathbb Q)\simeq\mathbb Q^r.
\]
Alexander duality in the oriented six-sphere gives
\[
 b_1(X\setminus A,\mathbb Q)=r.
\]
On the other hand, $f:X\setminus A\to U:=\mathbb P^1\setminus\Sigma$
is a proper smooth fibre bundle.  The base $U$ has the homotopy
type of a graph with $s-1$ loops.  The degree-one Leray--Serre
sequence therefore gives an exact sequence
\[
 0\longrightarrow H^1(U,\mathbb Q)
 \longrightarrow H^1(X\setminus A,\mathbb Q)
 \longrightarrow H^1(F,\mathbb Q)^{\pi_1(U)}
 \longrightarrow0.
\]
In particular, the correct formula retains the invariant
subspace:
\[
 r=s-1+\dim H^1(F,\mathbb Q)^{\pi_1(U)}\le s.
\]
Every nonempty fibre over $\Sigma$ has at least one irreducible
component, so $r\ge s$.  Equality follows throughout.  Thus
every fibre over $\Sigma$ has precisely one reduced irreducible
component and the invariant subspace is all of the
one-dimensional $H^1(F,\mathbb Q)$.  Fibres outside $\Sigma$ are
smooth and connected, hence irreducible.  This proves both
assertions without assuming monodromy triviality at the outset.
It is the rank-one repair of the component-count argument;
the unrestricted formula in \citet[Lemma~4.2]{CDP2020D} is not
being used.

Let $S=f^{-1}(t)_{\mathrm{red}}$.  A horizontal prime $H$ maps
surjectively to the base by properness, so $H\cap S\ne\varnothing$.
Since $S$ is irreducible, $S\subset H$ would imply $S=H$,
contrary to horizontality.  The canonical section of
$\mathcal O_X(H)$ consequently restricts to a nonzero section
on the reduced irreducible space $S$, with a nonempty zero
set.  Its zero divisor is a nonzero effective Cartier divisor
on $S$.  Indeed the section cannot vanish identically on
any local irreducible branch: its pullback would then vanish
on an open subset of the irreducible normalization, hence
everywhere by analytic continuation.  In the reduced local
ring it consequently avoids every minimal prime and is a
non-zero-divisor, including at self-intersections.

Let $\nu:S^\nu\to S$ be the normalization.  If $S^\nu$ were
in $\mathcal C$, it would have a compact K\"ahler resolution
$\mu:\widehat S\to S^\nu$; this precise surface statement is
\citet[Lemma~3.3 and Corollary~3.4]{Fujino2021E}.
The pulled-back line bundle
\[
 L=(\nu\mu)^*(\mathcal O_X(H)|_S)
\]
has $c_1(L)=0$ integrally, since $H^2(X,\mathbb Z)=0$.
Its pulled-back section is nonzero and has zeros: normalization
is finite surjective, and modification is proper surjective.
Thus it defines a nonzero effective divisor $B$ on
$\widehat S$.  For a K\"ahler form $\omega$,
\[
 0=\langle c_1(L)\smile[\omega],[\widehat S]\rangle
   =\int_B\omega>0,
\]
a contradiction.  The argument applies to the reduced support
even when the original fibre is nonreduced.  Moishezon spaces
are in $\mathcal C$, so the stated consequence follows.

For the assertion about a reducible fibre, let $S_i$ be one
of its reduced irreducible components.  The fibre is connected,
so another component $S_j$ meets $S_i$; otherwise $S_i$ would
be an open-and-closed component of its support.  Distinctness
ensures $S_i\not\subset S_j$.  Apply exactly the preceding
canonical-section and positivity argument to
$\mathcal O_X(S_j)|_{S_i}$.  Its first Chern class is zero
because $H^2(X,\mathbb Z)=0$, whereas its effective Cartier
zero divisor is nonempty.  This proves the additional assertion
without any hypothesis on the general fibre's first Betti number.

For the application, choose a very general smooth fibre $F$
of $g_0$ meeting a horizontal component of $E_0$ and avoiding
its vertical components.  As in E.2.6, adjunction gives
$c_1(K_F)=0$ and $F$ is minimal.  The identity of E.2.5 now
restricts to
\[
 K_F^{-A}=\mathcal O_F((-D)E_0|_F),
 \qquad -D>0,
\]
with a nonzero effective divisor on the right.  If a positive
power $K_F^m$ had a nonzero section, multiplying a suitable
power of that section by the $m$th power of the displayed
anticanonical section would give a nonzero holomorphic
function with zeros on $F$.  Compactness forbids this.
Thus $\kappa(F)=-\infty$.  Moreover $F$ is not K\"ahler,
since its nonzero effective divisor $(-D)E_0|_F$ has zero
real Chern class.  The Enriques--Kodaira classification then
places $F$ in class VII and gives $b_1(F)=1$; see
\citet[Chapter~VI, Theorem~1.1 and Table~10, pp.~243--244]
{BHPV2004E}.  The preceding argument applies with $f=g_0$
and any horizontal prime component $H$ of $E_0$.
\end{proof}

The proof supplies a necessary restriction on an unknown
smooth specialization.  It neither classifies the singular
non-K\"ahler normalizations nor replaces the unresolved
non-normal-fibre step in the holomorphic-fibration argument
of \citet[\S7]{CDP2020D}.  That step is not used here:
E.2.9--E.2.10 prove the needed horizontal-divisor vanishing,
and E.2.21--E.2.24 complete the actual specialization and
deformation-chain arguments.

% Owner: classification_e.  The proof does not import CDP2020 Theorem2.2.
\subsubsection*{Removing the negative-slope horizontal fixed part}

\begin{proposition}[E.2.9 (Proved): no horizontal divisor when all fibres are irreducible]
Let $X$ be a compact complex threefold diffeomorphic to $S^6$,
and let $f:X\to\mathbb P^1$ be a surjective holomorphic map
with connected fibres.  If every reduced fibre is irreducible,
then $X$ has no horizontal prime divisor for $f$.
\end{proposition}
\begin{proof}
Suppose that $H$ is such a divisor.  It is Cartier because
$X$ is smooth.  Since $H$ is horizontal and every reduced
fibre $S$ is irreducible, $H$ meets $S$ and does not contain
it.  As in E.2.8, restriction of its canonical section is a
non-zero-divisor on every local branch of $S$.

We first record the elementary vanishing used below.  On a
compact reduced irreducible complex space $S$, let $D$ be a
nonzero effective Cartier divisor and $V$ a vector bundle.
Multiplication by the powers of the canonical section
identifies
\[
 H^0(S,V(-nD))\subset H^0(S,V),\qquad n\ge0,
\]
with a decreasing sequence of finite-dimensional vector
spaces.  A section in their intersection vanishes to every
order along $D$.  Krull intersection in a local ring at a
point of $D$ makes its germ zero.  After pullback to the
irreducible normalization, analytic continuation makes the
section zero everywhere.  The decreasing sequence therefore
stabilizes at zero.  In particular
\[
 H^0(S,V(-nD))=0\quad\hbox{for all sufficiently large }n.
\]
The same conclusion holds for a vector bundle on a thickened
fibre $mS$.  Indeed $S$ is Cartier in $X$, and the finite
filtration by powers of its ideal has successive quotients
\[
 V|_S\otimes\mathcal O_S(-jS),\qquad 0\le j<m.
\]
Apply the preceding reduced-space statement to these finitely
many vector bundles, with the same nonzero divisor $H|_S$,
and then use their exact sequences.  This also proves that
the multiplication maps used here remain injective on $mS$:
its only associated component is $S$ and $H$ contains none
of its local branches.

Let $B$ be any holomorphic vector bundle on $X$.  The map
$f$ is flat, since its local rings are torsion-free over the
discrete valuation rings of the smooth base.  Every analytic
fibre $F_t$ is a Cartier, Cohen--Macaulay, Gorenstein surface,
possibly nonreduced.  Its dualizing line is
\[
 \omega_{F_t}=K_X|_{F_t},
\]
where the constant one-dimensional factor from the target
has been trivialized.  Serre duality on this compact
Gorenstein surface gives
\[
 H^2(F_t,B(nH)|_{F_t})^*
 \simeq H^0\bigl(F_t,(B^*\otimes K_X)(-nH)|_{F_t}\bigr).
\]
The compact-space duality being used includes singular
Cohen--Macaulay spaces; it is the absolute duality theorem
of \citet[\S3, Theorem~1, p.~81, and the complete-intersection
calculation on p.~91]{RR70push}, already specified in E.3.1,
not a smooth-fibre replacement.
The preceding filtration argument makes the right hand
side zero for all sufficiently large $n$, for each $t$.

For fixed $n$, the nonvanishing locus
\[
 Z_n=\{t\in\mathbb P^1:
        h^2(F_t,B(nH)|_{F_t})>0\}
\]
is closed by semicontinuity.  Multiplication on the dual
section spaces gives $Z_{n+1}\subset Z_n$.  Their intersection
is empty by fibrewise vanishing.  Compactness of $\mathbb P^1$
therefore supplies a single $N$ with $Z_N=\varnothing$.
Constant-zero cohomology and analytic cohomology-and-base-change
\cite[Chapter~III, Theorem~4.12(ii)]{BS76push} give
\[
 R^2f_*(B(NH))=0.
\]
Higher direct images vanish in degrees greater than two by
the fibre dimension.  The coherent sheaf $R^1f_*(B(NH))$
has vanishing first cohomology after a sufficiently positive
twist on $\mathbb P^1$.  The projection formula and Leray
sequence consequently give, for some integer $a\gg0$,
\[
 H^2(X,B\otimes M_+)=H^3(X,B\otimes M_+)=0,
 \qquad
 M_+=\mathcal O_X(NH)\otimes f^*\mathcal O(a).
\]
All coherence, semicontinuity and constant-zero base-change
uses are the proper analytic theorems audited in E.3.1;
the argument does not infer base change from coherence alone.

Independently, the reduced-space vanishing argument applied
to the connected smooth manifold $X$ and its nonzero divisor
$H$ gives
\[
 H^0(X,B\otimes M_-)=0,
 \qquad M_-=\mathcal O_X(-bH),\quad b\gg0.
\]
We must arrange these different vanishings for one line
bundle.  The exponential sequence gives
\[
 \operatorname{Pic}(X)\simeq H^1(X,\mathcal O_X)
\]
because $H^1(X,\mathbb Z)=H^2(X,\mathbb Z)=0$.
Let $\ell_-,\ell_+$ be the corresponding additive classes.
The affine line
\[
 \ell(z)=(1-z)\ell_-+z\ell_+,\qquad z\in\mathbb C,
\]
is realized by a holomorphic family of line bundles on
$X\times\mathbb C$: choose additive holomorphic Cech
cocycles representing the two classes and exponentiate their
linear combination.  Equivalently, use the integrable
Dolbeault operators obtained by a holomorphic linear
combination of fixed closed $(0,1)$-form representatives.
Thus no representability or universal-Picard-space assertion
is needed here.

For the proper projection $X\times\mathbb C\to\mathbb C$,
the cohomology jumping loci
\[
 T_j=\{z:h^j(X,B\otimes M_z)>0\},\qquad j=0,2,3,
\]
are closed analytic subsets of $\mathbb C$ by the proper
analytic finite-complex theorem
\cite[Chapter~III, Theorem~4.1]{BS76push}: the condition is
expressed by the rank minors in a local finite free complex.
In particular no closedness claim from the corrected
Grauert Section~7, Satz~2 is used.  The set $T_0$ omits $z=0$,
while $T_2$ and $T_3$ omit $z=1$.  Each is therefore proper,
and their finite union cannot fill the irreducible complex
line.  Choose $z$ outside that union.  Then
\[
 \chi(X,B\otimes M_z)=-h^1(X,B\otimes M_z)\le0.
\]
Every line bundle on $X$ is topologically trivial, so
Hirzebruch--Riemann--Roch makes this Euler characteristic
equal to $\chi(X,B)$.

Finally take $B=T_X$.  Since $H^2(X,\mathbb Z)=H^4(X,\mathbb Z)=0$,
the only contributing term in its Riemann--Roch formula is
\[
 \chi(X,T_X)=\frac12\langle c_3(X),[X]\rangle
            =\frac12\chi_{\mathrm{top}}(S^6)=1.
\]
This contradicts the preceding inequality.  Hence no
horizontal prime $H$ exists.
\end{proof}

\begin{proposition}[E.2.10 (Proved): persistence for either canonical slope]
Under the hypotheses of E.2.5, for every nonzero
$D=kl-3k-4l$ one has $E_0=0$ and $e=1$.
After shrinking the disc, the evaluation of
$\pi_*\mathcal L$ gives a proper flat holomorphic map
\[
 g:\mathcal X\longrightarrow\mathbb P^1\times\Delta
\]
whose restriction to every fibre is its intrinsic algebraic
reduction.  A smooth general central fibre of $g$ is a
complex two-torus.  Its normalized punctured period has
an interior limit as in E.2.7.

Two families satisfying E.2.5 cannot have the same smooth
central manifold unless their ordered multiplicity pairs
agree, regardless of the signs of their canonical slopes.
Uniqueness of a smooth completion at an interior period and
the arbitrary deformation-chain statement E.2.4 are not
asserted.
\end{proposition}
\begin{proof}
The case $D>0$ is E.2.6.  If $D<0$ and $E_0$ had a
horizontal prime component, the final part of E.2.8 would
give $b_1(F)=1$ for a smooth general fibre of $g_0$ and
would make every reduced fibre of $g_0$ irreducible.
E.2.9 would then forbid precisely that horizontal component.
Thus for either sign $E_0$ is vertical.

The connected-divisor-limit argument in the middle of the
proof of E.2.6 only uses verticality, not the sign of $D$.
Its generic choice of pencil section has central zero
support equal to the disjoint union of the vertical fixed
part and $e$ distinct regular fibres.  The total divisor
is reduced and irreducible with connected central support
after normalization and Stein factorization.  Consequently
$E_0=0$ and $e=1$.  Surjectivity of evaluation, flatness and
the torus assertion follow from the remainder of that same
proof.  The hypothesis of E.2.7 is now satisfied for either
sign, and gives the interior period limit.

Suppose two such one-sided families have the same central
manifold $X_0$.  Its algebraic reduction is intrinsic, so
both pencil lines are the same line
$L=g_0^*\mathcal O(1)$ up to isomorphism.  With indices
$i=1,2$, the two root identities give
\[
 K_{X_0}^{A_i}=L^{D_i},\qquad
 L^{D_1A_2-D_2A_1}=\mathcal O.
\]
The line $L$ is not torsion, because a positive power has
a nonzero section with zeros.  Hence $D_1/A_1=D_2/A_2$.
E.2.3 recovers the ordered coprime multiplicity pair from
this rational number.  This proves the common-limit
assertion without identifying the limit with a member of
the given period construction.
\end{proof}

The vanishing argument is an application of standard proper
cohomology, duality and Riemann--Roch machinery; no novelty
claim is intended.  It proves the needed absence of a
horizontal fixed divisor under the rank-one hypothesis.
It does not establish the different claim that every
holomorphic Hopf-surface fibration on a complex six-sphere
is impossible.

\subsubsection*{An extension criterion at an isolated base point}

\begin{proposition}[E.2.11 (Proved): purity with a smooth torus Jacobian]
Let $B\subset\C^2$ be a sufficiently small bidisc about $0$, put
$U=B\setminus\{0\}$, and let $g:X\to B$ be a proper flat
holomorphic map from a connected smooth complex $(n+2)$-fold,
with connected fibres of dimension $n$.  Let $G\to B$ be a
smooth proper holomorphic family of complex $n$-tori with a
holomorphic zero section and its group law.  Suppose that
$X_U\to U$ is a holomorphic $G_U$-torsor.  Then $g$ is a
submersion on all of $X$, and the given action extends uniquely
to make $X\to B$ a holomorphic $G$-torsor.

No polarization, K\"ahler metric on $X$, or section of the
given torsor on $U$ is assumed.  The smooth extension $G$ is
an explicit hypothesis.
\end{proposition}
\begin{proof}
Shrink $B$ so that the holomorphic Lie bundle of $G$ has a
frame $e_1,\ldots,e_n$.  Its period local system has a basis
of holomorphic sections
\[
 \lambda_1(b),\ldots,\lambda_{2n}(b)\in\C^n,
 \qquad
 G_b=\C^n/\sum_{a=1}^{2n}\Z\lambda_a(b).
\]
Indeed the lattice is locally constant over the simply
connected bidisc, and the local inverses of the relative
exponential map make its sections holomorphic.  The periods
are a real basis at every $b$.

The infinitesimal torsor action gives holomorphic vector
fields $V_1,\ldots,V_n$ on $X_U$.  Flatness makes
$Z=g^{-1}(0)$ have complex codimension two in $X$.
Hartogs extension in local coordinates therefore extends
each coefficient of each $V_j$ uniquely across $Z$; uniqueness
glues these local extensions to holomorphic vector fields
on $X$.  This use of Hartogs applies to functions on the
complement of an analytic set of codimension at least two
in a complex manifold, hence also coefficientwise to any
locally free sheaf.  It does not assert extension across a
divisor.  The identities
\[
 dg(V_j)=0,\qquad [V_i,V_j]=0
\]
hold on the dense open subset $X_U$ and consequently on $X$.

Every $V_j$, and every linear combination with coefficients
pulled back from $B$, is complete.  Its integral curve stays
in a compact fibre; a finite endpoint of a maximal solution
would have an accumulation point in that fibre, where the
holomorphic ordinary differential equation extends the
solution.  The flows thus define a holomorphic action of
$\C^n$ on $X$ over $B$.

For $a=(a_1,\ldots,a_{2n})\in\R^{2n}$ take the time-one
flow of the vertical holomorphic vector field whose Lie
coordinate is
\[
 \sum_{j=1}^{2n}a_j\lambda_j(g(x)).
\]
These maps depend smoothly on $a$ and add under composition:
all their coefficients are constant along every vertical
flow.  For $a\in\Z^{2n}$ the map is the identity on $X_U$,
hence everywhere by holomorphic uniqueness.  We obtain an
action of the fixed compact real torus
$T=\R^{2n}/\Z^{2n}$ by holomorphic automorphisms of $X$,
over $B$.  On the fibre over $b$, it is the holomorphic
action of $G_b$ obtained by factoring the $\C^n$ action by
its period lattice.  The action of $T$ is free on $X_U$.
The variable period identities also make the holomorphic
Lie-bundle action descend to a holomorphic action
$G\times_BX\to X$.  Its restriction therefore acts on
each scheme-theoretic fibre, including its nilpotent
structure.

We show that its stabilizers have dimension zero everywhere.
Let $x\in X$, $b=g(x)$, and let $H$ be the identity component
of the stabilizer.  Viewed in $G_b$, it is a compact complex
subgroup: it is the identity component of the inverse image
of $x$ under the holomorphic orbit map $G_b\to X_b$.
Put
\[
 K_x=\ker(dg_x)=T_x(X_b),
\]
where the right-hand side is the Zariski tangent space of
the scheme-theoretic fibre.  Differentiating the holomorphic
$G_b$ action on this fibre gives a holomorphic representation
$H\to\operatorname{GL}(K_x)$.  It is trivial, because every
holomorphic function on the compact connected complex group
$H$ is constant.  The induced representation on
$T_xX/K_x$ is also trivial: this quotient injects into
$T_bB$ by $dg_x$, and every group element acts over $B$.
Averaging a Hermitian inner product over the compact group
$H$ gives an invariant complement to $K_x$.  The full
derivative representation on $T_xX$ is therefore trivial.

Here the required local linearization has a direct proof.
Choose holomorphic coordinates $z$ centred at $x$ and a
sufficiently small $H$-invariant neighbourhood on which all
$z(hy)$ are defined.  Such a neighbourhood exists by
compactness of $H$ and continuity of the action.  With
normalized Haar measure set
\[
 \Psi(y)=\int_H z(hy)\,dh.
\]
This map is holomorphic, invariant under $H$, and has
derivative the identity at $x$.  It is a holomorphic
coordinate system after shrinking.  Its invariance then
forces every $h\in H$ to be the identity near $x$, and hence
on connected $X$.  This contradicts freeness on $X_U$ unless
$H$ is trivial.  Thus every stabilizer is finite, and the
$n$ vector fields $V_j$ are linearly independent everywhere.

Their commuting flows give local holomorphic coordinates
$(z,s)\in\C^n\times\C^2$ in which they span the $z$
directions: choose a transverse complex surface and apply
the holomorphic inverse function theorem to the flow map.
In these coordinates $g$ has the form
\[
 g(z,s)=h(s).
\]
Equidimensionality implies that the germ $h^{-1}(b)$ is
zero-dimensional.  Consequently the germ
$h:(\C^2,0)\to(B,b)$ is finite; this is the local finiteness
criterion obtained by Weierstrass division for an isolated
fibre.  If $dg_x$ failed to have rank two, the determinant
of $dh$ would vanish at $s=0$.  It is not identically zero,
since the map is a submersion away from $Z$.  Its zero set
would therefore contain a complex curve through $0$.
Finiteness sends that curve to a curve of critical values.
This contradicts smoothness over $B\setminus\{0\}$.
Hence $g$ is a submersion everywhere.

The $G_b$ orbits now have the full dimension of every smooth
connected fibre and are compact.  Each fibre is a single
orbit.  Its finite stabilizer acts trivially on its tangent
space, because $G_b$ is abelian; it also acts trivially on
the base tangent quotient.  Averaging as above, now over
that finite group, again makes its full derivative trivial
and then makes the action locally trivial.  Freeness on
$X_U$ forces the stabilizer to be trivial.  Thus the extended
action is free and transitive on every fibre.  Local
holomorphic sections of the submersion identify $X$ with
$G$ over small base opens, which proves the torsor assertion.
The extension is unique because it agrees with the original
one on the dense open subset $X_U$.
\end{proof}

\begin{remark}[Scope and literature check]
This is a conditional extension criterion, not a general
purity theorem for proper flat maps.  In particular it does
not apply to B.4.15: that construction has vertical
threefolds over the bidisc origin and is not flat over the
bidisc.  Searches on 15 November 2026 for purity and
discriminants of complex-torus fibrations, Hartogs extension
of torus actions, and isolated critical values did not
locate the exact statement under these hypotheses without
an added K\"ahler assumption.  We give the argument above
and claim no priority.
\end{remark}

% Owner: classification_e. Root and boundary_d full-source referee PASS.
% Uses the independently refereed abstract purity lemma E.2.11.
\begin{proposition}[E.2.12 (Proved): no additional central critical values]
In the smooth-specialization setting of E.2.10, let
\[
 g:\mathcal X\longrightarrow\mathbb P^1\times\Delta
\]
be the extended intrinsic fibration, and let
$c_1(s),c_2(s),c_0(s)$ be the three labelled critical-value
sections of E.2.7.  The critical values of $g_0$ are exactly
the distinct points among $c_1(0),c_2(0),c_0(0)$.
In particular, no additional critical value occurs only
on the central target curve.

If the normalizing projective transformations $A_s$ of
E.2.7 have a degenerate limit, the central regular fibres
are isomorphic complex two-tori and $g_0$ has at most two
critical values.  If the constant image of the limiting
rank-one matrix is $0$ or $1$, this torus is the
corresponding finite-root abelian surface.  If that image
is an ordinary source point, it is the ordinary torus at
that point and need not be an abelian surface.  The cusp
cannot be the constant image.
\end{proposition}
\begin{proof}
E.2.7 gives $0<|z(0)|<R_{k,l}$.  A local logarithm
therefore extends the normalized parameter $c(s)$
holomorphically to the disc.  The critical image is closed
and contains the three punctured sections, so it contains
their three limiting points.  Suppose that a different
central critical value $b_0=(v_0,0)$ exists.  Take a small
bidisc $B$ about $b_0$ which avoids all three section
closures and all other central critical values.  Put
$U=B\setminus\{b_0\}$.  The restricted family over $U$
is a proper smooth family of complex two-tori.  We verify
that its Jacobian extends to a smooth holomorphic torus
group family over $B$, so that E.2.11 applies.

The normalized source coordinate is $t=A_s(v)$.
Represent $A_s$ by a matrix with holomorphic entries,
after clearing poles and removing a common power of
$s$, so its central matrix is nonzero.  If that matrix
is invertible, $t$ extends holomorphically near $b_0$ and
its value there is ordinary: a marked value would put
$b_0$ on the corresponding section closure.  The
ordinary period model with parameter $c(s)$ then gives
the desired smooth Jacobian over $B$.

Otherwise the central matrix has rank one.  Denote its
constant image by $t_*$ and its exceptional target point
by $v_*$.  For each marked source value different from
$t_*$, its preimage tends to $v_*$.  At least two of the
three marked values are different from $t_*$, so $v_*$
belongs to at least two labelled section closures.
Consequently $v_0\ne v_*$.  After shrinking $B$, the
function $t(s,v)$ tends uniformly to $t_*$ there.
For an ordinary $t_*$ the ordinary period model again
extends directly.

Suppose instead that $t_*=0$ or $1$, and use a local
coordinate $\xi$ at that marked point.  The divisor of
$\xi(t(s,v))$ on $B$ is supported on $s=0$, since $B$
avoids the corresponding critical-value graph.
Consequently
\[
 \xi(t(s,v))=s^n u(s,v),\qquad n>0,
\]
where $u$ is a nowhere-zero holomorphic function.
This also explains why no zero of the leading
coefficient may be overlooked: such a zero would put
the point on the labelled graph closure.
Choose a nearby $v_1$ for which $(v_1,0)$ is regular.
The actual family over the full $s$-disc at $v_1$ is
smooth, so its integral degree-one monodromy is the
identity.  The parameter $c(s)$ has no winding.  By
the marked comparison of E.2.7, this monodromy is
$A_1^n$ or $A_2^n$, respectively.  Their orders are
three and four, so $e\mid n$, where $e=3$ or $4$.
A holomorphic $e$th root of $u$ exists on the bidisc,
and
\[
 r(s,v)=s^{n/e}u(s,v)^{1/e}
\]
is an actual holomorphic root coordinate on $B$.
Pull back the smooth finite-root torus period model
by $(c(s),r(s,v))$.  This produces the required
smooth torus group family $G\to B$.  The larger
multiplicities $3k,4l$ do not change this Jacobian
argument: their additional deck transformations
are translations, and the linear Jacobian model
is the original order-three or order-four root model.

For completeness $t_*=\infty$ is impossible.  In the
cusp coordinate the same divisor argument gives
$q=s^n u$ with $n>0$.  The smooth family at $v_1$
would have monodromy $A_0^n$.  But
$A_0=I+N$, $N^2=0$, and
$N\widehat u=\widehat w\ne0$, so $A_0^n\ne I$.

It remains to identify $G|_U$ with the actual Jacobian,
rather than merely comparing individual tori.
The punctured bidisc $U$ is simply connected.  Its
integral first-homology local system therefore has
a global marking.  Over $s\ne0$ that marking agrees
with the reference period marking, after one constant
integral change of basis; the transition comparisons
of E.1.3 have identity linear homology action at
fixed marked parameter.  Both holomorphic period
subbundles extend on $U$, and they agree on its dense
open subset $s\ne0$.  They agree everywhere on $U$.
Equivalently, the two quotients of their holomorphic
Lie bundles by the same marked period lattice are
isomorphic torus group families on $U$.  This is a
Jacobian identification and makes no assertion that
the actual affine torsor has a section on $B$.

The total space is smooth and $g$ is proper and flat.
The only possible nonsmooth fibre over $B$ lies above
$b_0$, and its codimension in the total space is two.
All hypotheses of E.2.11 now hold.  That lemma makes
$g$ smooth also above $b_0$, a contradiction.  The
central critical-value set is therefore precisely
the set of labelled limits.

For the final assertions, in the rank-one case every
ordinary central target point different from $v_*$
has the same limiting marked torus: it is obtained
from $(c(0),t_*)$, or from the finite-root value
$(c(0),r=0)$.  Equality of the extended period
subbundles proves this assertion without assuming
separatedness for threefold completions.
All marked values except possibly $t_*$ have their
preimages tending to $v_*$.  Hence there are at most
two central critical values; if $t_*$ is ordinary,
there is exactly one.
At $t_*=0$ or $1$, the finite linear automorphism of
the root torus has complex eigenvalues $1$ and a
primitive third or fourth root of unity.  The
rational idempotents for these two distinct
cyclotomic factors split that torus, up to isogeny,
as a product of two elliptic curves.  It is therefore
an abelian surface.  This argument makes no
projectivity assertion for an ordinary limiting
torus, and supplies no polarization on the whole
central threefold.
\end{proof}

\begin{remark}[Scope of E.2.11--E.2.12]
Purity here is obtained from the explicit good
Jacobian near a hypothetical extra critical value.
Smooth total space and flatness alone were not used
as a general purity theorem.  These results do not
give uniqueness of an interior smooth completion,
do not exclude the remaining isotrivial central
fibrations, and do not prove persistence of the
canonical pencil on an arbitrary new deformation
branch.
\end{remark}

% Owner: classification_e. Root and boundary_d full-source referee PASS.
\begin{proposition}[E.2.13 (Proved): persistence of an isolated finite multiple fibre]
In the setting of E.2.10, suppose the central value of
one finite labelled critical section is distinct from
the limits of the other two labelled sections.  Write
its multiplicity as $m$, with $m=3k$ or $4l$, and its
linear monodromy order as $e=3$ or $4$, respectively.
Near that section, normalize the base change taking
an $m$th root of its defining equation.

The resulting total space is smooth, and its map to
the root bidisc is a proper smooth family of two-tori.
The deck group acts freely on the total space.
Consequently the actual central reduced finite fibre
is smooth, has exactly multiplicity $m$, and belongs
to the same smooth family of finite free affine
quotients as the punctured reduced finite fibres.

This applies to both finite fibres when the three
labelled critical values have distinct limits.  In
the two-critical-value collapsed-target case, it
applies to the isolated finite labelled fibre.
\end{proposition}
\begin{proof}
Choose target coordinates $(v,s)$ in a bidisc about
the section, and let $\Gamma$ be its graph.  Take a
local defining function $\xi$ for $\Gamma$.
On the smooth fourfold the divisor identity is
\[
 g^*\Gamma=m\mathcal D,
\]
where $\mathcal D$ is the prime closure of the
punctured reduced multiple-fibre family.  Indeed
this equality holds away from $s=0$.  The only
possible extra divisorial component would be
$X_0$, but $g_0$ is nonconstant, so $X_0$ is not
mapped into the point $\Gamma(0)$.  The closure
$\mathcal D$ is Cartier because the total space
is smooth.  Thus locally, even if its central
restriction is initially singular or reducible,
\[
 g^*\xi=a^m u
\]
for a local equation $a$ of $\mathcal D$ and a
holomorphic unit $u$.

Let $B'\to B$ be the smooth root bidisc
$r^m=\xi$, and let $Y$ be the normalization
of $\mathcal X_B\times_B B'$.  A local $m$th
root of $u$ factors its defining equation as
\[
 r^m-a^m u=
 \prod_{\zeta^m=1}(r-\zeta a u^{1/m}).
\]
Its normalized branches are separate copies of
the original smooth local fourfold.  Therefore
$Y\to\mathcal X_B$ is a finite etale cover,
and $Y$ is smooth.  The deck group permutes these
branches freely, including over $a=0$.
This is an assertion about the normalized cover;
the unnormalized fibre product generally has
intersecting branches.

The induced map $g':Y\to B'$ is proper and
equidimensional of relative dimension two.
Since both source and base are smooth, the
two pulled-back base parameters form a regular
sequence in each local ring, giving flatness.
Over $s\ne0$ this is the actual smooth root
torus family.  Its connectedness makes $Y$
connected and its generic fibre connected.
The Stein factor of $g'$ is finite birational
over the normal bidisc and hence is the
bidisc itself, so all fibres are connected.
Equivalently this follows from normality of
$Y$ and the finite birational form of its
Stein factorization.

By E.2.12 the original central map has no
critical value away from the labelled limits.
Our choice of the section therefore makes
$g'$ smooth away from the origin of $B'$.
It remains to supply a smooth Jacobian over
the full root bidisc.

If the normalizing projective transformations
$A_s$ have an invertible limit, the source
coordinate at the finite marked value has
the form
\[
 \xi_{\rm src}=\xi\,u_1(v,s)
\]
with $u_1$ a unit.  Thus its order-$e$ root
is $r^{m/e}u_1^{1/e}$, and the exact finite-root
period model with parameter $c(s)$ extends
smoothly to the full root bidisc.

If $A_s$ has rank-one limit, the isolated
section in the hypothesis must be the
preimage of its constant image $t_*$:
all other marked preimages tend to the
exceptional point.  Thus $t_*$ is precisely
the finite marked value under consideration.
Near the isolated section the source
coordinate has divisor equal to that section
plus a positive multiple of $s=0$, so
\[
 \xi_{\rm src}=s^n\xi\,u_1(v,s),\qquad n>0,
\]
with a nowhere-zero unit $u_1$.  Choose a
nearby central regular target point with
$\xi\ne0$.  Its actual smooth torus family
has identity parameter monodromy.  As in
E.2.12, $c(s)$ has no winding and hence
$e\mid n$.  The required holomorphic root
coordinate on $B'$ is now
\[
 s^{n/e}r^{m/e}u_1^{1/e}.
\]
It again gives a smooth reference Jacobian
over $B'$.  On the punctured bidisc its
marking and holomorphic period subbundle
agree with the actual Jacobian by the same
simply-connected argument as E.2.12.

E.2.11 makes $g'$ a smooth torsor over the
entire root bidisc.  In particular $r$ is
a submersion coordinate on the central
threefold upstairs.  The unramified deck
cover then shows that the original fibre
has exactly multiplicity $m$.  Restricting
to $r=0$ gives a proper smooth family of
two-tori in the parameter $s$, with free
affine deck action.  Its quotient is a
proper smooth family of the reduced finite
fibres, proving all assertions.
\end{proof}

\begin{remark}[Scope of the finite-chart argument]
This argument identifies the smooth root charts
at the isolated finite critical values.  The
excluded target-collision cases are treated in
E.2.14--E.2.15.  The actual cusp is recovered
separately in E.2.16--E.2.20, and E.2.21 glues
these comparisons into the full one-sided
completion theorem.
\end{remark}

% Owner: classification_e. Root and boundary_d full-source referee PASS.
\begin{proposition}[E.2.14 (Proved): extension of the circle field and exclusion of finite-root collapse]
In the setting of E.2.10, the normalized vector field
$\delta$ of the punctured family extends to a holomorphic
relative vector field on $\mathcal X$ after shrinking
the parameter disc.  Its time-one flow is the identity,
and its real circle action on the central fibre is
equivariantly diffeomorphic to the prescribed circle
action on the punctured fibres.  In particular its
central fixed locus is a smooth rational curve with
normal bundle $\mathcal O(-1)\oplus\mathcal O(-1)$
and the same two opposite circle weights.

The rank-one limit of the normalizing projective
transformations $A_s$ cannot have constant image
$0$ or $1$.  More generally, a collapsed constant
image is impossible whenever the limiting torus's
complex $\delta$ line has two real-linearly
independent periods.  Thus the finite-root
isotrivial possibilities in E.2.12 are excluded.
\end{proposition}
\begin{proof}
Let $\pi:\mathcal X\to\Delta$ be the original
proper holomorphic submersion.  The coherent sheaf
\[
 \mathcal H=\pi_*T_{\mathcal X/\Delta}
\]
is torsion-free: multiplication by a nonzero base
function is injective on the vector bundle upstairs
and hence on its sections.  Its generic rank is
one by E.2.1.  On the smooth disc it is therefore
locally free of rank one.  Choose a generator $V$
after shrinking the disc.  On $\Delta^*$ it has
the form
\[
 V|_{X_s}=b(s)\delta_s,
\]
with $b$ holomorphic and nowhere zero there.
The normalized fields $\delta_s$ agree on overlaps
of local classifying maps: at fixed discrete data
the comparison maps of E.1.3 have fibre-linear
part the identity and fix the marked period
$\widehat\delta$.  Their translations also
preserve the invariant vector field.
Moreover $V$ is vertical for $g$, because this
holds on the dense punctured family.

Choose a regular central target value and a small
bidisc around it on which $g$ is smooth.  The
ordinary or finite-root comparison used in
E.2.12 supplies its smooth marked Jacobian, with
the marked period $\widehat\delta$ extending as
a nonzero holomorphic Lie section.  The resulting
invariant field $\delta_B$ on the actual torsor
is nonvanishing there.  Comparing one nonzero
coefficient of $\delta_B$ with $V$ shows that
$b(s)$ extends holomorphically to $s=0$.
Its independence of the target coordinate follows
from that equality on $s\ne0$ and holomorphic
uniqueness.

If $b(0)=0$, the restriction $V|_{X_0}$ vanishes
on a nonempty open subset of connected $X_0$,
and therefore vanishes everywhere.  In local
coordinates for the submersion $\pi$, every
coefficient of $V$ is then divisible by $s$.
The quotients glue uniquely to a global
holomorphic relative vector field $V/s$.
This contradicts the fact that $V$ generates
the free module $\mathcal H$.  Hence $b(0)\ne0$,
and $V/b(s)$ gives the required global extension
of $\delta$.

The extended field is complete on each compact
fibre of $\pi$.  Its time-one flow is the
identity on the punctured family and hence on
the whole family by holomorphic uniqueness.
Its real flow consequently defines an action
of $\mathbb R/\mathbb Z$ on $\mathcal X$
over the parameter disc.  Average a smooth
metric and an Ehresmann horizontal distribution
over this compact group.  Parallel transport
then gives equivariant diffeomorphisms of
the fibres.  The central action is therefore
the same effective semifree action, with a
single fixed two-sphere, as on the prescribed
fibres.

The fixed set of a holomorphic compact-group
action is a complex submanifold, as follows
from the averaging coordinates used in
E.2.11.  Here it is a proper smooth family
of curves of genus zero, and hence each fibre
is $\mathbb P^1$.  Its two holomorphic normal
weight lines have the same degrees as in
the punctured family.  Those degrees are
$-1,-1$ by C.3.11, or directly by the two
fixed toric edge normals.  The characters
distinguish the lines, so their central
restrictions are exactly
$\mathcal O(-1)\oplus\mathcal O(-1)$.

Suppose now that $A_s$ has rank-one limit
with constant image a finite marked value.
E.2.12 identifies every central regular
torus, with its marked $\delta$ direction,
with the corresponding finite-root torus
at parameter $c(0)$.  The invariant lattice
of $A_j$ has rank two and spans its
one-dimensional complex fixed subspace.
It contains $\widehat\delta$; its second
primitive period is independent over
$\mathbb R$.  The complex fixed subspace
is precisely the $\delta$ line, since
the other complex eigenvalue is a
primitive third or fourth root of unity.
Consequently, in the normalization in
which the $\delta$ period is $1$, there
is a constant second period
\[
 \kappa\in\mathbb C\setminus\mathbb R.
\]
This is a period of the invariant field
on every central regular fibre.  The
time-$\kappa$ flow is therefore the
identity on their dense union, and
hence on all of $X_0$.

The flow of the nonzero field $\delta_0$
thus factors through a nontrivial
holomorphic action of the compact
elliptic curve
$E=\mathbb C/(\mathbb Z+\kappa\mathbb Z)$.
Such an action has no fixed point on
a connected complex manifold unless
it is trivial.  Indeed its derivative
at a fixed point is a holomorphic
representation of the compact complex
group into a general linear group,
so is constant.  Averaging local
coordinates over the compact group
then makes the action trivial near
that point, and holomorphic uniqueness
makes it trivial everywhere.

The elliptic action here is nontrivial
because $\delta_0\ne0$.  It therefore
has no fixed point.  This contradicts
the fixed rational curve of the
extended real circle action: a zero
of the generating field is fixed
by its entire complex flow.  Equivalently,
the resulting nowhere-zero real vector
field would give $\chi(X_0)=0$,
whereas $X_0$ is diffeomorphic to $S^6$.
The same argument applies to an ordinary
constant-image torus if its $\delta$
line has a second independent period.
\end{proof}

\begin{remark}[What the elliptic-action argument does not assume]
The argument uses the global one-dimensional
space of vector fields on the known punctured
members and the actual marked period comparison
on a central regular bidisc.  It does not
assume that all translations of an ordinary
central torus extend over the singular fibre.
For an ordinary torus the $\delta$ line may
have only the one specified real period;
such a case is not excluded by the
elliptic-action argument alone.
\end{remark}

\begin{proposition}[E.2.15; Proved; an isotrivial volume obstruction]
Let $f:X\to\PP^1$ be a proper holomorphic map with connected
fibres from a compact connected complex threefold.  Suppose
that for some $p\in\PP^1$ the restriction over
$U=\PP^1\setminus\{p\}$ is a smooth family of mutually
isomorphic complex two-tori.  If
\[
 K_X^{\otimes A}\simeq f^*\mathcal O_{\PP^1}(D),
 \qquad A\in\Z_{>0},\quad D\in\Z,
 \tag{E.2.15.1}
\]
then necessarily $D<-A$.  In particular the ordinary-source
collapsed-target case of E.2.12 cannot occur for the
six-sphere families under consideration.
\end{proposition}
\begin{proof}
Choose an affine coordinate $t$ on $U\simeq\C$ and put
$q=t^{-1}$ near $p$.  Since $U$ is simply connected, its
integral first-homology local system has a global marking.
The holomorphic period subbundle is constant: locally the
marked periods belong to the countable integral-change-of-basis
orbit of one fixed torus, and a continuous map from a connected
base into such a countable subset of a complex Grassmannian
is constant.  For the last assertion, any two distinct image
points can be separated by a real continuous coordinate;
its connected image would contain an interval and would be
uncountable.  Thus the Jacobian group family is the product
with one fixed torus $T$.

Choose a nonzero translation-invariant holomorphic two-form
$\eta_T$ on $T$.  Local holomorphic origins identify our
family with its Jacobian, and two choices differ by a
translation.  Hence $\eta_T$ defines a global nonvanishing
relative canonical form $\eta$ on $X_U$.  Its fibre volume
is a fixed positive number:
\[
 \int_{X_t}\frac14\eta\wedge\overline\eta=V>0.
 \tag{E.2.15.2}
\]
This construction uses neither a global torsor section nor
a polarization on $X$.  The canonical determinant
isomorphism for a submersion gives the nonvanishing
absolute three-form $\eta\wedge dt$ on $X_U$; this notation
does not require choosing an absolute lift of $\eta$.

Use (E.2.15.1) to pull back a meromorphic section of
$\mathcal O(D)$ whose divisor is $D[p]$.  We obtain a
meromorphic pluricanonical form $\sigma$ with
\[
 \operatorname{div}(\sigma)=D f^*(p).
\]
On $X_U$ it is holomorphic and nowhere zero.  Its quotient
by $(\eta\wedge dt)^{\otimes A}$ is a holomorphic function
constant on each compact connected fibre.  It therefore
descends to an entire nowhere-zero function $\phi$ on $\C$:
\[
 \sigma=\phi(t)(\eta\wedge dt)^{\otimes A}.
 \tag{E.2.15.3}
\]

On a small disc about $p$, the meromorphic section
$q^{-D}\sigma$ is in fact a holomorphic nowhere-zero section
of $K_X^{\otimes A}$ over the whole inverse image.  This
follows directly from (E.2.15.1), or by cancellation of the
actual divisors, with their multiplicities.  Denote it by
$\Omega_A$.  On the punctured inverse image, since
$dt=-q^{-2}dq$, it is
\[
 \Omega_A=(-1)^A P(q)(\eta\wedge dq)^{\otimes A},
 \qquad P(q)=q^{-D-2A}\phi(1/q).
 \tag{E.2.15.4}
\]
The function $P$ is single-valued and holomorphic on the
punctured disc.  No extension of a comparison with another
threefold is asserted here.

A nowhere-zero holomorphic $A$-canonical form determines
a smooth positive density: if it is
$h(z)(dz_1\wedge dz_2\wedge dz_3)^{\otimes A}$, take
$|h|^{2/A}$ times the ordinary complex-coordinate volume
density.  This is independent of the coordinates.  The
density associated with $\Omega_A$ has finite integral
over the inverse image of a smaller closed disc, because
that inverse image is compact.  Fibre integration of
(E.2.15.4), using (E.2.15.2), consequently gives
\[
 \int_{0<|q|<\epsilon}|P(q)|^{2/A}\,d\mathrm{Area}(q)
 <\infty.
 \tag{E.2.15.5}
\]
Only a fixed positive normalization constant is suppressed.
This step concerns a smooth positive density on the actual
total space and requires no K\"ahler metric.

We spell out the isolated-singularity consequence of
(E.2.15.5), including when $2/A<1$.  The function
$|P|^{2/A}$ is subharmonic.  If $|q|=r$, the mean-value
inequality on the disc of radius $r/2$ centred at $q$ gives
\[
 |P(q)|^{2/A}
 \leq \frac{4}{\pi r^2}
       \int_{r/2<|z|<3r/2}|P(z)|^{2/A}\,d\mathrm{Area}(z).
\]
The last integral tends to zero as $r\to0$, by
integrability.  Hence $r^A|P(q)|\to0$, uniformly on
$|q|=r$.  The function $q^A P(q)$ therefore has a removable
singularity and vanishes at zero.  For some holomorphic
$H$ on the disc,
\[
 P(q)=q^{1-A}H(q).
 \tag{E.2.15.6}
\]
Equivalently, the pole order of $P$ is at most $A-1$.

Combining (E.2.15.4) and (E.2.15.6) yields
\[
 \phi(1/q)=q^{D+A+1}H(q).
 \tag{E.2.15.7}
\]
If $D\geq-A$, the right-hand side extends to zero at
$q=0$.  Thus $\phi$ extends holomorphically to $\PP^1$
and vanishes at $p$.  A holomorphic function on $\PP^1$
is constant, so $\phi=0$, contradicting its construction.
This proves $D<-A$.

In E.2.10 the actual canonical-root equality has
\[
 A=12kl,\qquad D=kl-3k-4l,
 \qquad
 \frac DA=\frac1{12}-\frac1{4l}-\frac1{3k}
       \geq-\frac12>-1.
\]
If the collapsed constant source point is ordinary,
E.2.12 says that the central map has precisely one
critical value and is isotrivial on its complement.
All hypotheses above apply to the actual central
threefold, and give the required contradiction.
\end{proof}

\begin{remark}[Scope of the volume argument]
The proof uses the equality of line bundles in
(E.2.15.1), not just the topological vanishing of
$c_1(K_X)$.  It also uses constancy of the torus periods,
which makes the fibre volume in (E.2.15.2) constant.
It does not assert a general good-reduction or
comparison theorem for non-K\"ahler torus degenerations.
Together with the separate elliptic-action exclusion
in E.2.14, it removes the two types of collapsed
target permitted by E.2.12.  Identification of a
noncollapsed interior completion, and the behaviour
on arbitrary new deformation branches, remain
distinct questions.
\end{remark}

% Owner: classification_e. Root and boundary_d full-source referee PASS.
\begin{proposition}[E.2.16 (Proved): the central cusp is reduced and irreducible]
In the one-sided smooth specialization of E.2.10,
the three labelled critical values remain distinct.
The two finite fibres have their original exact
multiplicities $m=3k$ and $n=4l$.  Every reduced
central critical fibre is irreducible, and the
cusp fibre has multiplicity one.  Its underlying
surface $S$ satisfies
\[
 \chi_{\rm top}(S)=2,\qquad b_1(S;\mathbb Q)=2.
\]
This supplies the reduced irreducible Cartier central
surface used in the analytic cusp identification of
E.2.17--E.2.20.
\end{proposition}
\begin{proof}
E.2.12, E.2.14 and E.2.15 exclude every degenerate
limit of the normalizing projective transformations.
They therefore extend as invertible projective
transformations, and the three labelled critical
values remain distinct.  E.2.13 proves the stated
finite-fibre multiplicities.

Normalize the three sections to $0,1,\infty$.
Over $B^\circ=\mathbb P^1\setminus\{0,1,\infty\}$,
the relative map is a proper submersion over
$B^\circ\times\Delta$.  Hence its integral
first-homology monodromies on the central fibre
are the same matrices $A_1,A_2,A_0$ as in
E.2.7.  The common invariant space of the dual
matrices $A_1^t,A_2^t$ has dimension one:
in the displayed marking it is generated by
the covector $\gamma$.  This is obtained by
solving their simultaneous fixed equations;
the exact calculation is recorded in
\path{code/e_push_interior_cusp.py}.

Let $r$ be the total number of irreducible
components of the three reduced critical
fibres, and let $X^\circ$ be their complement
in $X_0$.  Alexander duality on $S^6$ gives
\[
 b_1(X^\circ;\mathbb Q)=r.
\]
The degree-one Leray sequence over the
punctured projective line, whose second
cohomology vanishes, gives instead
\[
 b_1(X^\circ;\mathbb Q)
 =b_1(B^\circ;\mathbb Q)
   +\dim H^1(T;\mathbb Q)^{\rm mon}
 =2+1=3.
\]
There are three nonempty critical fibres,
so each has exactly one reduced component.

Write the cusp fibre as $dS$ with $d\ge1$.
We prove $d=1$ by a direct meridian argument.
The smooth fibration over $B^\circ$ has
connected fibres, hence
$\pi_1(X^\circ)\to\pi_1(B^\circ)$ is
surjective.  The inclusion
$X^\circ\hookrightarrow X_0$ is surjective
on fundamental groups, and its kernel is
normally generated by meridians about the
three reduced divisor components.  To see
the latter assertion, make a null-homotopy
disc transverse to their smooth loci and
disjoint from their singular loci, which
have real codimension at least four.
Removing small discs about the finitely
many intersections expresses its boundary
as a product of conjugate meridians.

A meridian of a multiplicity-$a$ component
maps to the $a$th power of the corresponding
base meridian, since locally the base
function is a unit times the $a$th power
of a reduced local equation.  Consequently
there is a surjection
\[
 \pi_1(X_0)\longrightarrow
 \left\langle a,b,c\ \middle|\
 a^m=b^n=c^d=abc=1\right\rangle.
 \tag{E.2.16.1}
\]
This is the usual orbifold fundamental
group, but the surjection has just been
proved without assuming a stack
factorization.  Its source is trivial.

For clarity the group on the right is
nontrivial for every $d\ge2$ here.
Except for $(m,n,d)=(3,4,2)$, one has
\[
 \frac1m+\frac1n+\frac1d<1.
\]
Indeed for $d\ge3$ the sum is at most
$11/12$; for $d=2$, increasing either
$m=3k$ or $n=4l$ above $3,4$ makes it
at most $23/24$.  Put
$\alpha=\pi/m$, $\beta=\pi/n$,
$\theta=\pi/d$, and choose $\ell>0$ with
\[
 \cosh\ell=
 \frac{\cos\theta+\cos\alpha\cos\beta}
      {\sin\alpha\sin\beta}>1.
\]
Let
\[
 R_\alpha=\begin{pmatrix}
 \cos\alpha&\sin\alpha\\
 -\sin\alpha&\cos\alpha
 \end{pmatrix},\qquad
 D_\ell=\operatorname{diag}(e^{\ell/2},e^{-\ell/2}).
\]
In $\operatorname{PSL}_2(\mathbb R)$ the
elements represented by $R_\alpha$ and
$D_\ell R_\beta D_\ell^{-1}$ have orders
$m,n$, and their product has trace
$-2\cos\theta$ and projective order $d$.
Taking its inverse for $c$ gives a
nontrivial representation of (E.2.16.1).
For the exceptional triple, the
permutations $a=(123)$ and $b=(1324)$
have orders three and four, and
$ab=(24)$ has order two, giving a
nontrivial representation in $S_4$.
Thus $d\ge2$ contradicts (E.2.16.1),
and the cusp is reduced.

The ordinary torus fibres have Euler
characteristic zero.  The finite reduced
fibres are smooth finite free quotients
of two-tori by E.2.13, so also have
Euler characteristic zero.  Additivity
over the three singular values therefore
gives $\chi_{\rm top}(S)=\chi(X_0)=2$.
Finally, the valid irreducible-fibre
degree-one invariant-cycle statement
\cite[Corollary 3.7, arXiv version~1]{HondaViaclovsky2025D}
gives
\[
 b_1(S;\mathbb Q)=
 \dim H^1(T;\mathbb Q)^{A_0}=2.
\]
Its irreducibility hypothesis has just
been proved.  This use is separate from
the false unrestricted equality
converse audited in B.4.17.
\end{proof}

% Owner: root. classification_e and boundary_d full-source referee PASS.
\begin{proposition}[E.2.17 (Proved): persistence of the normal-crossing cusp]
In the smooth fixed-pair specialization of E.2.10,
the cusp divisor is a relative normal-crossing divisor
over the parameter disc.  More precisely, after
shrinking that disc, at every point of the total cusp
divisor there are holomorphic coordinates
$(z_1,z_2,z_3,s)$, with $s$ the original parameter,
in which the cusp coordinate of the intrinsic map is
one of
\[
 q=z_1,\qquad q=z_1z_2,\qquad q=z_1z_2z_3.
 \tag{E.2.17.1}
\]
There are exactly two relative triple-point sections.
Both lie on the fixed rational curve of E.2.14.
In particular the normalization of the total cusp
divisor is smooth and proper over the parameter
disc, and its central fibre is the normalization
of the central cusp.
\end{proposition}
\begin{proof}
By E.2.14--E.2.16 the three labelled critical
values remain distinct, the cusp has multiplicity
one, and its reduced central fibre is irreducible.
We normalize its target section to $q=0$.
On a sufficiently small target disc, zero is the
only critical value on every parameter fibre.
The extended field $\delta$ is zero precisely on
the relative fixed rational curve $\mathcal C$,
and the two normal characters there are $+1,-1$.
For $s\ne0$ the two triple points of the prescribed
cusp belong to $\mathcal C_s$; outside that curve
the cusp is smooth or has an ordinary double crossing.

We first record exactly the local singularity result
we use.  If $h_s(a,b)$ is a holomorphic family and
$h_0=0$ is a reduced plane-curve germ with an
isolated critical point, choose a sufficiently
small representative of that germ.  The sum of
the multiplicities of all nearby critical points
of $h_s$, regardless of their critical values,
is $\mu(h_0)$.  One can see this directly from
the analytic critical scheme
\[
 \mathbb C\{a,b,s\}/
       ((h_s)_a,(h_s)_b).
\]
After taking a sufficiently small representative,
it is finite over the parameter disc.  Its two
equations are a regular sequence, and its
central fibre is zero-dimensional.  It is thus
Cohen--Macaulay of dimension one, with $s$ a
nonzerodivisor, and is finite flat over the
disc.  Its degree is the central Jacobian-algebra
dimension, namely $\mu(h_0)$.

The non-splitting theorem of
\citet[Theorem 2 and its corollary, pp.~274--275]{Lazzeri1973Root}
says that if a deformation of an isolated
hypersurface singularity has more than one
singular point in a single deformed fibre,
their summed Milnor numbers are strictly
smaller than the original Milnor number.
The strict numerical inequality is the
conclusion of the cited corollary's proof.
Consequently, if all the critical points
just counted have critical value zero,
there can be only one such point.  This
application uses a flat deformation of
plane hypersurfaces: the map
$(a,b,s)\mapsto(h_s(a,b),s)$ is flat near
the central germ because its two base
parameters form a regular sequence in
the smooth source.  No assertion of
separatedness of compact threefolds enters
this local argument.

Let $p$ be a central cusp point outside
$\mathcal C_0$.  A relative holomorphic
flow box for the nonzero field $\delta$
gives coordinates $(a,b,v,s)$ in which
$\delta=\partial_v$.  Invariance of the
intrinsic map gives
\[
 q=h_s(a,b).
\]
The central equation is reduced, by
E.2.16.  A reduced plane-curve germ is
either smooth or has an isolated critical
point.  In the smooth case the implicit
function theorem already gives the first
form in (E.2.17.1).

In the singular case every critical point
of $h_s$ in a small representative lifts
to a critical point of the actual map
on $X_s$.  All its critical values are
therefore zero.  For $s\ne0$ each such
point is a node, since the flow box is
disjoint from the fixed curve and hence
from the two triple points.  Conservation
and non-splitting give one node and
$\mu(h_0)=1$.  Thus the central Hessian
is nondegenerate.  The implicit function
theorem gives its unique critical-point
section, whose critical value is zero.
The holomorphic Morse lemma with
parameters then gives $q=ab$ throughout
this neighbourhood, with $s$ unchanged.

It remains to treat the fixed curve.
Holomorphic coordinates obtained by
averaging over the extended compact
circle may be chosen in the form
$(z,x,y,s)$, with $\mathcal C$ defined
by $x=y=0$ and with circle weights
$0,+1,-1,0$.  The original parameter
$s$ is retained by this averaging.
The holomorphic invariant function $q$
vanishes on $\mathcal C$, so its
convergent power series has the form
\[
 q=F_s(T,z)=T g_s(T,z),\qquad T=xy.
 \tag{E.2.17.2}
\]
The coefficient
$Q_s(z)=g_s(0,z)$ is intrinsically a
section of the dual of the product of
the two normal weight lines along
$\mathcal C_s$.  By E.2.14 this is
$\mathcal O_{\mathbb P^1}(2)$.

The central section $Q_0$ is not
identically zero.  Otherwise the
central equation in every such chart
would be divisible by $x^2y^2$,
contradicting reducedness of the
central cusp.  Its zero divisor
therefore has degree two.  Suppose
that both zeros coincide at a point
$z=0$.  The plane equation
$F_0(T,z)=Tg_0(T,z)$ is reduced:
a repeated factor of $g_0$ would
remain repeated after $T=xy$, and
$T$ does not divide $g_0$ because
$Q_0$ is not identically zero.
Thus $F_0$ has an isolated plane
critical point there.

For $s\ne0$ the two zeros of $Q_s$
are the two distinct triple points
of the prescribed cusp.  Under the
supposed collision they both lie
in the chosen small plane
representative.  They are distinct
critical points of $F_s$ with value
zero.  Every other critical point
of $F_s$ also has value zero:
this is automatic at $T=0$, and
at $T\ne0$ it follows by lifting
to $xy=T$, where the quotient
map $(x,y,z)\mapsto(T,z)$ is a
submersion.  All sufficiently small
$T$ have such a lift in the original
coordinate chart.  A nonzero
critical value would contradict
the actual critical-value set of
the intrinsic map.

Conservation of all critical
multiplicities and Lazzeri's
non-splitting theorem now give a
contradiction to the two distinct
critical points.  Hence $Q_0$ has
two simple zeros.  They extend
as two disjoint holomorphic
sections over the parameter disc.
Where $Q_0\ne0$, the factor
$g_s(xy,z)$ is a unit and can
be absorbed into $y$, giving
$q=xy$.  At a simple zero,
$\partial g_0/\partial z\ne0$;
replace $z$ by $g_s(xy,z)$ to
obtain $q=xyz$, again keeping
$s$ unchanged.  This proves
all local forms in (E.2.17.1).

Properness makes the central cusp
compact, so finitely many such
neighbourhoods cover it and a
single shrinking of the parameter
disc suffices.  The normalization
of a relative normal-crossing
divisor separates its local
branches.  In the coordinates
above every normalized branch
is smooth over $s$, and these
descriptions agree on overlaps.
Normalization is finite for
complex analytic spaces; its
composition with the proper
cusp map is therefore proper.
Its fibres are precisely the
normalizations of the cusp
fibres in these local charts.
This proves the final assertions.
\end{proof}

\begin{remark}[Literature and remaining identification]
Non-splitting is a known theorem of
Lazzeri, not a new singularity-theoretic
claim.  We opened the original paper
and use its strict numerical
corollary.  The application combines
that theorem with the actual extended
circle action, reducedness, and the
absence of extra critical values
proved above.  Searches for
non-splitting, Milnor-number
conservation and degenerations of
torus-fibred complex six-spheres
did not locate this application;
no priority is claimed.  The
proposition supplies relative
normal-crossing charts and a
simultaneous normalization.  It
does not by itself identify the entire
analytic cusp neighbourhood.  The
prescribed marked Mumford completion
is recovered by E.2.18--E.2.20.
\end{remark}

% Owner: classification_e. Root and boundary_d full-source referee PASS.
\begin{proposition}[E.2.18 (Proved): the normalized cusp and its marked hexagon]
Let $\nu:\mathcal S\to\mathcal W$ be the
simultaneous normalization of the relative cusp
in E.2.17, and let $\mathcal D\subset\mathcal S$
be its reduced conductor.  After shrinking the
parameter disc, each fibre pair is the standard
toric pair
\[
 (\mathcal S_s,\mathcal D_s)
 \simeq(\mathrm{dP}_6,\hbox{its toric hexagon}).
\]
The six conductor components may be labelled over
the disc so that their incidence, self-intersections
$-1$, opposite-pair involution, and endpoint
identifications are those of the punctured family.
The relative logarithmic tangent bundle has a
global holomorphic frame:
\[
 T_{\mathcal S/\Delta}(-\log\mathcal D)
 \simeq\mathcal O_{\mathcal S}^{\oplus2}.
\]
This assertion identifies the normalized pair
and supplies a relative frame; it does not by
itself identify the analytic smoothing neighbourhood.
\end{proposition}
\begin{proof}
In the relative charts of E.2.17, the normalized
branches are smooth over the original parameter.
Their conductor branches and their intersections
are also smooth over that parameter.  The
normalization of the conductor is consequently
a proper smooth family of curves, and the
normalized intersection strata are finite
etale over the disc.  Their connected components
and incidence are constant.  Since the disc
is simply connected, label the six curve
families $\mathcal D_1,\ldots,\mathcal D_6$
and the six vertex sections by continuation
from a punctured fibre.  Each curve fibre
is a smooth curve of genus zero.  Its two
distinct vertex sections remain distinct;
the same relative local charts exclude a
new self-intersection or an additional
incidence at the centre.  Thus the central
conductor is a hexagon of six embedded
rational curves.  The degree of each
relative normal line is constant in its
proper curve family, so all six central
self-intersections equal $-1$.

Write $S=\mathcal S_0$ and $D_i=\mathcal D_{i,0}$.
Ehresmann transport gives $b_1(S)=0$ and
$b_2(S)=4$.  Contract the three disjoint
curves $D_1,D_3,D_5$.  The analytic
Castelnuovo contraction gives a smooth
compact complex surface $Y$; the needed
analytic statement, including smoothness
of the resulting point, is
\citet[\S4, no.~8(a), pp.~364--365]{Grauert1962E}.
Only these three contractions on the
central surface are used; no relative
contraction theorem is assumed.

The images $C_2,C_4,C_6$ of the remaining
curves are smooth rational curves.  Each
meets two contracted curves transversely
once, so has self-intersection
$-1+1+1=1$.  Each pair meets once, at
the image of their intervening contracted
curve; these three intersection points
are distinct and the intersections are
transverse.  Thus they form a triangle.
Also $b_1(Y)=0$ and $b_2(Y)=1$.

Any one of these curves, say $C=C_2$,
has nonnegative intersection with every
effective divisor: distinct irreducible
curves have nonnegative local intersection
multiplicities and $C^2=1$.  Adjunction
gives $K_Y\cdot C=-3$.  A nonzero section
of $K_Y^{\otimes r}$ for $r>0$ would
give an effective divisor of intersection
$-3r$ with $C$, impossible.  Hence
$\kappa(Y)=-\infty$.  The compact-surface
classification
\cite[Chapter~VI, Theorem~1.1 and Table~10]{BHPV2004E}
together with $b_1(Y)=0$ makes $Y$ rational:
class~VII has first Betti number one,
and a ruled surface has the first Betti
number of its base curve.  A rational
surface with $b_2=1$ is $\mathbb P^2$.
The curves $C_2,C_4,C_6$ are lines,
since their self-intersections are one.
Their three pairwise intersection points
are distinct, so they form a projective
triangle.  Blowing up its three vertices
recovers exactly the toric degree-six
del Pezzo surface and its hexagonal
boundary.  This proves the assertion
about the central pair; the punctured
pairs were already the prescribed ones.

The conductor involution is holomorphic
over the disc.  Its permutation of the
six labelled curve families, and its
permutation of the vertex sections,
are locally constant.  They are therefore
the same opposite-pair and endpoint
identifications as on a punctured fibre.
The remaining scalar gluing parameters
are not claimed to be determined by
this discrete observation.

Finally put
$\mathcal E=T_{\mathcal S/\Delta}(-\log\mathcal D)$.
For every fibre pair just identified,
the two toric infinitesimal actions
trivialize its logarithmic tangent
bundle.  Thus
\[
 \mathcal E|_{\mathcal S_s}\simeq
 \mathcal O_{\mathcal S_s}^{\oplus2},
 \qquad h^0(\mathcal S_s,\mathcal E_s)=2.
\]
Constant-cohomology base change, in the
form used in E.3.1, makes the direct
image a rank-two holomorphic vector
bundle on the disc and identifies its
fibres with these section spaces.
Choose a frame after shrinking the
disc.  The evaluation map from its
pullback to $\mathcal E$ is an
isomorphism on every fibre, hence
is an isomorphism everywhere by
Nakayama's lemma.  This gives the
claimed relative frame without
choosing a potentially singular
family of projective coordinate
changes on the punctured disc.
\end{proof}

% Owner: boundary_d. Full proof submitted for root/E independent referee.
\begin{proposition}[E.2.19 (Proved): extension of the marked cusp torus action]
Let
\[
 f:\mathcal N\longrightarrow
 B=\Delta_q\times\Delta_s
\]
be the proper cusp neighbourhood of the smooth
specialization in E.2.17, and put
$\mathcal W=f^{-1}(q=0)$.  After shrinking the
bidisc, the rank-two vertical logarithmic tangent
bundle is trivial and is generated by two commuting
complete holomorphic vector fields.  Their action
factors through a holomorphic action of
$(\C^*)^2$ over $B$, with the integral marking
given by the periods $\mathbf w,\boldsymbol\delta$.
On the normalization of every $\mathcal W_s$
this is its usual toric action.

The marked rank-two lattice is primitive in the
period lattice of each ordinary torus fibre and
equals the full invariant lattice of the cusp
monodromy.  At the two relative triple-point
sections, the three nonzero tangent characters
are constant in $s$ and equal to those of the
prescribed Mumford cusp.  In the toric chart
lattice coordinates, transported from that
punctured model by their fixed integral change
of basis from the period marking, they are
respectively
\[
 \{(1,0),(0,1),(-1,-1)\},\qquad
 \{(-1,0),(0,-1),(1,1)\}.
 \tag{E.2.19.1}
\]
The assertion constructs the action on the actual
neighbourhood.  It does not assert that an
arbitrary equivariant quotient is smooth or
that the entire completion has already been
identified with a prescribed toric quotient.
\end{proposition}
\begin{proof}
\emph{The logarithmic bundle and its restriction.}
The coordinates in E.2.17 make $\mathcal W$ a
relative normal-crossing divisor over $s$.
Define
\[
 \mathcal L=
 \ker\left[
 d\log q:
 T_{\mathcal N/\Delta_s}(-\log\mathcal W)
       \longrightarrow\mathcal O_{\mathcal N}
 \right].
 \tag{E.2.19.2}
\]
The map is surjective: if
$q=z_1\cdots z_k$, it sends each
$z_i\partial_{z_i}$, $i\leq k$, to one.
Consequently $\mathcal L$ is locally free of
rank two.  Its sections are actual holomorphic
vector fields, tangent to the branches, and
annihilate both $q$ and $s$.

Let $\nu:\mathcal S\to\mathcal W$ and
$\mathcal D\subset\mathcal S$ be the
simultaneous normalization and conductor.
Restriction to a branch gives a canonical
isomorphism
\[
 \nu^*(\mathcal L|_{\mathcal W})
 \simeq T_{\mathcal S/\Delta_s}(-\log\mathcal D).
 \tag{E.2.19.3}
\]
This can be checked at every stratum.  For
$q=xy$, a basis is
$x\partial_x-y\partial_y,\partial_z$;
on the branch $x=0$ it becomes
$-y\partial_y,\partial_z$, a logarithmic
basis there.  For $q=xyz$, use
$x\partial_x-z\partial_z$ and
$y\partial_y-z\partial_z$.  Their
restrictions to $x=0$ have coefficient
matrix of determinant one in the basis
$y\partial_y,z\partial_z$.
The one-branch case is immediate.
The same computation keeps $s$ fixed.

By E.2.18 the right-hand side of
(E.2.19.3) has a global holomorphic
frame.  On the two preimages of each
conductor curve, the identifications
in (E.2.19.3) describe descent by an
invertible two-by-two matrix.  Its
entries are holomorphic on a proper
$\mathbb P^1$ family over $\Delta_s$,
so are constant along each curve
fibre and are holomorphic functions
of $s$.  Denote these matrices by
$G_j(s)$.

For $s\ne0$ both invariant logarithmic
fields in the prescribed Mumford model
descend.  Indeed, its period maps in
multiplicative coordinates multiply
the two coordinates by units depending
only on the base.  Their derivative
on the two vertical logarithmic
fields is the identity.  The toric
small-resolution charts preserve
these same fields.  Thus the entire
two-dimensional space of sections
on the normalized pair descends for
$s\ne0$.  In any chosen frame this
means $G_j(s)=I$, and the identity
theorem gives $G_j(0)=I$ as well.

There is no additional descent
condition hidden at a triple point.
For a normal-crossing local ring,
functions on the normalization
descend precisely when their
restrictions to every pair of
branches agree.  This is the usual
equalizer for
$\C\{x,y,z\}/(xyz)$, and likewise
for one or two branches.  Tensoring
with the locally free bundle
$\mathcal L|_{\mathcal W}$ proves
the same assertion for its sections.
Consequently the chosen normalized
frame descends, and
\[
 \mathcal L|_{\mathcal W}
       \simeq\mathcal O_{\mathcal W}^{\oplus2}.
 \tag{E.2.19.4}
\]

\emph{A frame on the whole neighbourhood.}
The map $f$ is proper and flat; this
also follows locally from the equations
$(q,s)=(z_1\cdots z_k,s)$.
The bundle $\mathcal L$ is flat over
$B$.  On each ordinary fibre it is
the tangent bundle of a complex
two-torus.  On each $q=0$ fibre it
is trivial by (E.2.19.4).  All these
fibres are compact and connected,
so in every case
\[
 h^0(\mathcal N_{q,s},
       \mathcal L|_{\mathcal N_{q,s}})=2,
\]
and those sections generate the
bundle pointwise.  The proper
constant-cohomology base-change
theorem, in the form established
in E.3.1, makes $f_*\mathcal L$
a rank-two vector bundle on $B$.
Its evaluation map
\[
 f^*f_*\mathcal L\longrightarrow\mathcal L
 \tag{E.2.19.5}
\]
is an isomorphism on every fibre,
and hence everywhere by Nakayama's
lemma.  Shrink the bidisc and
choose a holomorphic frame $V_1,V_2$.

These fields commute.  On every
ordinary compact torus, every
holomorphic vector field is a
translation field, so their
bracket vanishes there.  The
ordinary locus is dense, and
the bracket is holomorphic.
Each field is complete because
it is vertical and every fibre
of the proper map is compact:
the holomorphic local flow can
be continued for every complex
time on that compact invariant
fibre.  Holomorphic dependence
on initial points and parameters
gives a holomorphic action
\[
 \operatorname{Act}:\C^2\times
       \mathcal N\longrightarrow\mathcal N
\]
over $B$.

\emph{Holomorphic extension of the two toric periods.}
Choose a smooth point $p_0$ of the
central cusp away from its conductor.
The map $f$ is a submersion there,
so there is a holomorphic section
$p(q,s)$ through $p_0$.
On the central normalized toric
pair the two fields are a basis
of the toric logarithmic Lie
algebra.  Its open orbit is
$(\C^*)^2$.  Let
$v_1^0,v_2^0\in\C^2$ be a
basis of its full period lattice
in the coordinates $V_1,V_2$.
They are complex-linearly
independent.

Near each $v_i^0$, write
$\operatorname{Act}(v,p(q,s))=p(q,s)$
using two vertical coordinates
in a submersion chart around
$p_0$.  The derivative in $v$
is invertible at $(v_i^0,0,0)$,
since the infinitesimal action
is an isomorphism on that open
orbit.  The holomorphic implicit
function theorem gives a unique
holomorphic vector $v_i(q,s)$
near $v_i^0$ solving this equation.
For $q\ne0$, fixing $p(q,s)$
means that this translation is
the identity on the entire
torus fibre.  The equality of
holomorphic maps then extends
over $q=0$.  Thus $v_i(q,s)$
is a period on every fibre.
Their complex-linear
independence persists after
shrinking $B$.

Use these two columns as a
change of the frame.  Because
their coefficients are pulled
back from $B$, the new fields
still commute and are complete.
They have period one, and the
action consequently factors
through
$\C^2/\Z^2\simeq(\C^*)^2$,
using the exponential
$z\mapsto e^{2\pi i z}$
in each coordinate.

\emph{Primitivity and the actual marking.}
On $B^\circ=\Delta_q^*\times\Delta_s$
the period lattices of the
translation action are the
integral first-homology local
system of the torus family.
The continuously varying
periods $v_1,v_2$ represent
flat integral classes: in a
local differentiable
trivialization their flow
loops vary continuously, and
integral homology is discrete.
They generate a rank-two
invariant sublattice $L$.

Suppose $L$ were not primitive
at some ordinary fibre.  There
would be a rational pair
$(a_1,a_2)\notin\Z^2$ such
that $a_1v_1+a_2v_2$ is an
integral period there.  This
is an equality of flat
integral classes after
multiplication by a common
denominator, so it persists
throughout the connected
ordinary base.  In particular
the corresponding fixed
finite-order element of
$(\C^*)^2$ acts identically
there.  By the identity
theorem it acts identically
on the central open orbit.
This contradicts the choice
of $v_1^0,v_2^0$ as the full
central toric period lattice.
Thus $L$ is primitive.

The marked cusp monodromy,
constant in $s$ by E.2.16,
has the form
\[
 \widehat\gamma\longmapsto
       \widehat\gamma-\widehat\delta,
 \qquad
 \widehat u\longmapsto\widehat u+\widehat w,
 \qquad
 \widehat w\longmapsto\widehat w,
 \qquad
 \widehat\delta\longmapsto\widehat\delta.
 \tag{E.2.19.6}
\]
Its full integral invariant
lattice is
$\Z\widehat w\oplus\Z\widehat\delta$.
Since $L$ is a primitive
rank-two sublattice of this
rank-two lattice, equality
follows.  An integral
unimodular change of the two
periods therefore gives
exactly the marking
$\mathbf w,\boldsymbol\delta$.
This also proves that no
extra finite kernel has been
silently introduced.

\emph{Characters at the triple points.}
The two triple-point sections
are fixed by the torus action:
on the normalization they
are toric vertices.  Restrict
the derivative representation
to the compact real torus
$(S^1)^2$ and average a
Hermitian metric along either
section.  Its characters
are integral pairs and
vary continuously in $s$;
hence their multiset is
locally constant.  Equivalently,
Fourier averaging gives the
individual character
subbundles.  They coincide
at $s=0$ with the known
characters for $s\ne0$.

For clarity, the two unit
triangles in the prescribed
cocharacter lattice have
height-one rays
\[
 \{(0,0,1),(1,0,1),(0,1,1)\},
 \qquad
 \{(1,0,1),(0,1,1),(1,1,1)\}.
\]
Inverting their unimodular
ray matrices gives the two
character triples in
(E.2.19.1); the sum in
each triple is zero, as
required by invariance of
$q=z_1z_2z_3$.
The calculation is recorded
in \path{code/d_push_cusp_torus_action.py}.
All action and marking
assertions follow.
\end{proof}

% Owner: root. Working proof independently reviewed by classification_e and boundary_d.
\begin{proposition}[E.2.20 (Proved): identification of the actual analytic cusp]
In the smooth specialization of E.2.10, normalize
the three distinct critical-value sections to
$0,1,\infty$ and write $q=1/t$ near the cusp.
Let $c(s)$ be the holomorphic interior period
lift supplied by E.2.7.  After shrinking the
$(q,s)$-bidisc, its inverse image in the actual
family is biholomorphic, over that bidisc, to
the prescribed marked Mumford cusp family at
parameter $c(s)$.  The isomorphism induces
the given identification of the regular
torus Jacobians and their integral markings.
Its affine origin may be fixed by choosing
a holomorphic section through a smooth
point of the reduced cusp.
\end{proposition}
\begin{proof}
Write $f:\mathcal N\to\Delta_q\times\Delta_s$
for the proper actual cusp neighbourhood and
$\mathcal W=f^{-1}(q=0)$.  E.2.16--E.2.18
give its reduced relative normal-crossing
divisor, its normalized toric hexagon, and
the unchanged opposite-side identifications.
E.2.19 gives an actual holomorphic action
of $H=(\C^*)^2$ over the bidisc.  Its
period marking is the primitive vanishing
lattice
\[
 \Lambda_{\rm tor}
   =\Z\widehat w\oplus\Z\widehat\delta.
 \tag{E.2.20.1}
\]
We reconstruct the neighbourhood from
this action, keeping track of the
extension of its characters.

\emph{The topology and the exact covering.}
The relative coordinates
$q=x,xy,xyz$ give smooth real lifts
of the two real parameter directions
of $s$ which annihilate $dq$ and
are tangent to every cusp stratum.
Away from the cusp such lifts exist
because $f$ is a submersion.
A smooth partition of unity preserves
both these linear requirements.
On the inverse image of a smaller
closed $q$-disc the lifts integrate
along a sufficiently short parameter
path: $q$ remains constant, and
properness prevents escape.
Thus their flows identify the cusp
neighbourhoods, all strata, and
the regular-fibre inclusion maps
by $q$-preserving diffeomorphisms.

For the known Mumford cusp we
check the covering calculation
directly.  Each maximal toric
chart above $|q|<\epsilon$ is
\[
 \{(x,y,z)\in\C^3:|xyz|<\epsilon\},
\]
which is star-shaped and simply
connected.  Every intersection
contains the same connected dense
open torus $(\C^*)^2\times\Delta_q^*$.
Each connected component of an
intersection is open and must meet
that dense set, so the intersection
is connected.  The same holds for
the intersection with a finite
union of preceding charts.
Van Kampen proves simple
connectedness of each finite
union; a loop lies in such a
union by compactness.  The full
toric cover is therefore simply
connected.  The known free,
properly discontinuous period
action has deck group
\[
 \overline\Lambda=\Lambda/\Lambda_{\rm tor}
                 \simeq\Z^2.
\]
On a regular fibre its covering
is $\C^2/\Lambda_{\rm tor}\to
\C^2/\Lambda$.  The regular-fibre
map to the cusp body's fundamental
group is exactly the projection
to $\overline\Lambda$, and the
$q$-meridian dies through the
section $(x,y,z)=(q,1,1)$.
These statements also agree
with the original manuscript's
Proposition~4.7 and Corollary~4.8,
whose original PDF was checked
\cite[p.~30]{AC26Root}; the
preceding argument verifies
the step we use.  The
$q$-preserving transport transfers
this exact description to the
actual neighbourhood.

Take the universal covering
\[
 \widetilde{\mathcal N}\longrightarrow\mathcal N.
\]
Its deck group is
$\overline\Lambda$, and its
central normalized components
form the planar unfolding of
the hexagon.  The $H$ action
lifts.  Indeed the fundamental
loops of $H$ at a central open
orbit lie in one normalized
del Pezzo component, which is
simply connected, and hence die
in $\pi_1(\mathcal N)$.
The identity lift is unique
and satisfies the action law.
Choose a holomorphic section
$p(q,s)$ through a smooth
central cusp point and lift
it to the cover.  On $q\ne0$
its orbit map is an isomorphism
\[
 H\times(\Delta_q^*\times\Delta_s)
       \xrightarrow{\ \sim\ }
 \widetilde{\mathcal N}|_{q\ne0}.
 \tag{E.2.20.2}
\]
The stabilizer has been removed
exactly by (E.2.20.1), using
the just-verified regular-fibre
inclusion map.  In particular
the two characters of $H$
are globally defined on this
open covering space.

Connectedness of $H$ makes
its lifted action commute
with every deck transformation:
conjugation acts continuously
on the discrete deck group
and is the identity at the
identity of $H$.  On
(E.2.20.2), the two deck
generators therefore act by
\[
 (h,q,s)\longmapsto(P_i(q,s)h,q,s),
 \qquad i=1,2.
 \tag{E.2.20.3}
\]
The functions $P_i$ are the
actual marked period multipliers.
Changing the affine origin
does not change them.
The marked period comparison
of E.2.7 and E.2.12 identifies
them with the prescribed
$q$-monomials times holomorphic
units at $c(s)$.

\emph{Character extension across adjacent components.}
For a central component $D_v$,
choose a holomorphic base
section $a_v(q,s)$ through
its dense $H$ orbit.  Its
orbit map is an open
holomorphic embedding
\[
 H\times(\Delta_q\times\Delta_s)
       \longrightarrow\widetilde{\mathcal N}.
\]
It is injective on the
central open orbit and on
each regular fibre, and
its differential is invertible
along those loci.  Its
image contains the whole
regular covering.  Two
such origins differ for
$q\ne0$ by a holomorphic
translation $h_{vw}(q,s)$
in $H$.  We must prove
that this function has no
essential singularity at
$q=0$.

If both origins lie in
the same central component,
their ratio extends as a
holomorphic $H$-valued
unit.  The central open
orbit is free and transitive,
so this follows from the
holomorphic inverse-function
theorem for the orbit map,
with both parameters retained.

For two adjacent components,
use a crossing chart with
$q=ab$ and an orbit coordinate
$z\ne0$.  The stabilizer
of a general conductor
point is a primitive
connected one-dimensional
subtorus of $H$.  Its
two normal characters are
$+1,-1$: on each normalized
toric component these are
the primitive boundary
characters, and invariance
of $q$ makes them opposite.
Compact-circle averaging
gives these exact local
linear coordinates.  A
complementary torus direction
provides the orbit coordinate;
an invariant unit in $q$
can be absorbed while
preserving the weights.

Fix small nonzero constants
$a_0,b_0,z_0$ in this chart
and take the two special
origins
\[
 (a,b,z)=(q/b_0,b_0,z_0),\qquad
 (a,b,z)=(a_0,q/a_0,z_0).
 \tag{E.2.20.4}
\]
The stabilizer cocharacter
evaluated at $a_0b_0/q$
sends the first to the
second.  This equality
is justified by the local
action itself: a path
in that subgroup can
increase the modulus of
$a$ and decrease that
of $b$, with $ab=q$
fixed, staying inside
the crossing chart;
the phase can be varied
without changing either
modulus.  Thus the
possibly large group
parameter causes no
unjustified use of a
local coordinate outside
its domain.

Comparing (E.2.20.4) with
arbitrary chosen origins
adds only holomorphic
$H$-valued units, by the
same-component argument.
The adjacent comparison
therefore has the form
\[
 h_{vw}(q,s)=q^{n_{vw}}u_{vw}(q,s),
 \qquad n_{vw}\in\Z^2,
 \tag{E.2.20.5}
\]
where $u_{vw}$ is holomorphic
with values in $H$ on the
whole bidisc.  Every central
component is reached by
a finite adjacent path.
Multiplying these expressions
proves the same assertion
for its comparison with
the initial origin.
Different paths agree
on the regular cover,
and hence define identical
meromorphic characters.
The exponents are locally
constant in $s$, since
(E.2.20.5) holds with a
nonvanishing unit on the
full bidisc.

\emph{The fan and the analytic coordinate maps.}
At either triple-point
section the whole $H$
fixes the point.
Compact-torus averaging
gives three eigen-coordinates
with the constant integer
characters of E.2.19.
Their only primitive
positive invariant relation
is their product.
The reduced normal-crossing
equation therefore writes
$q$ as that product times
a holomorphic invariant
unit.  The unit is a
function of the product
and $s$; absorb it into
one eigen-coordinate to
obtain exactly $q=xyz$.

The orders of the two
global characters along
$D_v$, together with
$\operatorname{ord}_{D_v}q=1$,
give a ray $(v,1)$.
At a triple point the
three incident rays
form a unimodular cone.
Indeed all adjacent
orders and tangent
characters are constant
in $s$ and coincide
with the known ones
for $s\ne0$.
The entire resulting
fan is consequently
the prescribed $A_2$
triangulation, with
its fixed labels.
No additional refinement
is being introduced.

In a triple chart a
global character is
a Laurent monomial in
$x,y,z$ times a function
of $q,s$.  Evaluate it
on a section keeping
two nonzero coordinates
fixed and solving the
third from $q$.
The origin comparisons
above show that the
remaining function is
a power of $q$ times
a holomorphic unit.
Thus the dual monomial
coordinates of the
unimodular fan cone
pull back to $x,y,z$
times holomorphic units
in $q,s$.  They extend
holomorphically and
give a local
biholomorphism through
the triple point.
The crossing and
open-orbit charts give
the same conclusion
on the other strata.
The maps agree on
the dense regular
cover and therefore
glue.

We finally check the
common neighbourhood
and the deck action.
We do not take an
infimum of infinitely
many chart radii.
Choose finitely many
charts covering one
compact normalized
central del Pezzo
component and finitely
many overlaps for
the two deck generators.
Shrink the bidisc once
for this finite data.
On the regular cover
the character map
intertwines the actual
deck action with the
reference multipliers
$P_i(q,s)$ in
(E.2.20.3).  Both maps
extend in the charts
just constructed, so
their equality extends
to the central strata.
The reference multipliers
preserve the prescribed
invariant fan.  Deck
translates of the finite
chart collection cover
the entire central
unfolding and use the
same base bidisc.

We have obtained local
biholomorphisms at all
central strata, bijections
on their labelled torus
orbits, and the original
isomorphism on the whole
regular cover.  The
inverse coordinate maps
are the same unimodular
monomial charts with
the inverse units.
Quotienting by the
identical free properly
discontinuous deck
actions gives the
claimed biholomorphism
of actual analytic
cusp neighbourhoods.
\end{proof}

\begin{remark}[Scope and source comparison]
The toric covering and its
period action are the known
Mumford construction.  The
point here is to recover
that construction from
the actual smooth limit,
using the action proved
in E.2.19.  In particular
the absence of an essential
singularity in the character
comparison is proved in
(E.2.20.4)--(E.2.20.5);
it is not an assumed
meromorphic comparison.
The local identification
does not alone identify
the global threefold or
settle arbitrary deformation
chains.  E.2.21 supplies the
finite affine and global collar
comparison; E.2.22--E.2.24
provide the additional propagation
and closure arguments for general
bases and continuous structures.
\end{remark}

% Owner: classification_e. Root and boundary_d full-source referee PASS.
\begin{proposition}[E.2.21 (Proved): global identification of a smooth completion]
Let $\pi:\mathcal X\to\Delta$ satisfy the hypotheses of
E.2.5, with fixed coprime multiplicities $(3k,4l)$ on
every noncentral fibre.  Then $X_0$ is a member of that
same constructed family.  Its isomorphism parameter is
the interior value $z(0)$ of E.2.7.

After choosing a holomorphic lift $c(s)$ of that parameter
and shrinking the disc, the actual family is isomorphic
to the prescribed family at $c(s)$.  The ambiguity in
this comparison is the connected translation automorphism
group already determined in E.1.
\end{proposition}
\begin{proof}
E.2.10 extends the intrinsic pencil to a proper flat map
\[
 g:\mathcal X\longrightarrow\mathbb P^1\times\Delta.
\]
E.2.7 supplies its interior nonzero period $z(s)$ and hence
a holomorphic logarithmic lift $c(s)$.  E.2.12 excludes
additional critical values, and E.2.14--E.2.15 exclude
target collapse.  Normalize the three distinct labelled
sections to $0,1,\infty$ by a holomorphic projective
transformation over the disc.

On the ordinary locus the actual marked Jacobian has
period matrix $\Pi_{c(s)}(t)$.  More intrinsically, its
holomorphic period kernel in the complexified marked
lattice equals that of this matrix.  This equality holds
for $s\ne0$ by the given family identification and extends
over $s=0$ by the identity theorem.  The local marking
choices agree: an automorphism of a fixed marked member
has identity action on ordinary $H_1$, by E.1.2.  The
actual smooth family at any fixed ordinary target point
also has identity homology return around $s=0$.

E.2.13 gives smooth actual normalized root covers at the
two finite values, of degrees $3k$ and $4l$, with their
free cyclic deck actions.  Their marked Jacobians are
the corresponding reference root Jacobians.  E.2.20
identifies the actual cusp neighbourhood over the full
$(q,s)$-bidisc with the prescribed marked Mumford
neighbourhood, including its period multipliers.
It remains to identify the finite affine actions and
the global translations on the overlaps.

\emph{The finite affine classes.}
Choose local holomorphic origins on a finite root torus
family.  Every deck map is affine on the compact torus
fibres.  Its linear part is the fixed marked finite
monodromy.  Subtracting this part gives a holomorphic
translation section.  Lift it to the Lie bundle on
simply connected refinements of the root charts.
The finite-order relation says that the sum of its
successive linear translates is an integral period.
This lattice vector is locally constant in $s$:
the lifted translation and period vectors vary
continuously, whereas the integral coefficients are
discrete.  Changing the lift changes it by precisely
the integral coboundary allowed in E.1.

Consequently the finite affine cocycle has the same
integral class at $s=0$ as at $s\ne0$.  After matching
these classes, the remaining difference is a cocycle
in the complex vector sheaf.  The explicit finite
averaging in E.1 removes it, since division by the
finite group order is allowed in that sheaf.
This constructs the finite-root comparisons with
holomorphic dependence on $s$.  Equivalently one
may use a finite chart cover and its deck translates;
the same averaging solves the equivariant comparison.
The reduced finite fibres are the resulting free
quotients; they have not been assumed to be tori.

\emph{The integral collar classes.}
Use the two finite-root neighbourhoods, the cusp
neighbourhood, ordinary simply connected charts,
and fixed annular collars to cover the target.
Choose local holomorphic origins on the ordinary
torus charts.  A transition between two such
descriptions is fibrewise affine.  The common
integral marking makes its linear part the
prescribed one, leaving a translation section.

Lift these sections to the Lie bundles on simply
connected refinements of the collars.  Differences
on deck translates and sums on triple overlaps
are integral periods.  Every resulting integral
vector is constant in $s$.  This comparison takes
place on fixed collars away from the singular
centres; it does not evaluate a logarithm at a
singular fibre.  The finite and cusp chart
identifications specify exactly which integral
changes of origins are allowed.

For $s\ne0$ these classes are those of the
prescribed logarithmic-transform data, modulo
the integral changes of origins in E.1.
They are therefore the same at $s=0$.
One can normalize them once on the connected
parameter disc.  This includes the cusp
winding class, since E.2.20 identifies the
actual marked period action there.  Thus
there is no additional central integral
translation or fan datum.

\emph{Solving the holomorphic remainder.}
After this discrete normalization, the remaining
comparison is the additive cocycle in the
extension sheaf $\mathcal V$ of E.1.  Its
finite extension conditions are exactly
those of the already identified root
charts.  Its cusp extension condition
is the actual toric collar condition
proved in E.2.20.  The calculation in
E.1 therefore applies to this actual
cover and gives
\[
 0\longrightarrow\mathcal O_{\mathbb P^1}
       \boldsymbol\delta
 \longrightarrow\mathcal V
 \longrightarrow\mathcal O_{\mathbb P^1}(-\infty)
 \longrightarrow0,
 \qquad H^1(\mathbb P^1,\mathcal V)=0.
\]
The explicit finite averaging and Laurent-series
splitting in E.1 kill the additive cocycle.
The resulting translations glue the local
comparisons to a biholomorphism of $X_0$
with the prescribed member at $c(0)$.

The same solution works over a smaller parameter
disc.  On the fixed collars, Cauchy integral
formulas for the Laurent coefficients preserve
holomorphic dependence on $s$, as does finite
averaging.  Thus the solved translations and
their glued comparison are holomorphic in
$s$.  Their ambiguity is the connected
scalar translation group of E.1.  This
proves both the central and family assertions.
\end{proof}

\begin{remark}[Scope of E.2.21]
The known-pair locus is open by E.2.2 and
is now stable under full punctured-disc
smooth specializations.  These two facts
alone do not prove that an arbitrary open
known locus in a reduced analytic base
is topologically closed.  A holomorphic
disc through its boundary need not have
its entire punctured image in that locus.
The needed generic propagation argument
is proved in E.2.22.  E.2.23--E.2.24 then
establish the arbitrary-base and continuous-family
conclusions asserted in E.2.4.
\end{remark}

% Owner: classification_e. Root and boundary_d full-source referee PASS.
\begin{proposition}[E.2.22 (Proved): propagation from a known open subset of a disc]
Let $\pi:\mathcal X\to\Delta$ be a proper holomorphic
submersion with fibres diffeomorphic to $S^6$.  Suppose
there is a nonempty open subset $U\subset\Delta$ such
that every $X_s$, $s\in U$, is a member of the fixed
coprime multiplicity family $(3k,4l)$ of E.2.1.
Then every fibre over $\Delta$ is a member of that
same family.
\end{proposition}
\begin{proof}
Local miniversality in E.2.1 supplies actual
holomorphic comparisons with the prescribed family
locally over $U$, with locally constant integral
markings.  Thus the known open set has the relative
root and normal-crossing charts used below; only
pointwise isomorphisms are not being substituted
for these relative comparisons.

The proof first propagates the actual geometric data
outside finitely many successive discrete analytic
exceptional sets.  It then fills these sets in reverse
order by E.2.21.  Later exceptional sets are only
claimed to be discrete on the domain then under
consideration; they may accumulate at points removed
at an earlier stage.

\emph{The intrinsic line and the relative pencil.}
Put $A=12kl$ and $D=kl-3k-4l\ne0$.
The argument of E.2.5 constructs the unique line
bundle $\mathcal L$ on the whole total space with
\[
 \mathcal L^D=K_{\mathcal X/\Delta}^{A}.
\]
Indeed Ehresmann transport over the contractible disc
gives $H^1(\mathcal X,\Z)=H^2(\mathcal X,\Z)=0$,
so the exponential sequence identifies its Picard
group with a complex vector space.  This construction
needs no hypothesis on the fibres outside $U$.

Proper coherence makes $\mathcal V=\pi_*\mathcal L$
a torsion-free, hence locally free, sheaf on the
disc.  Its rank is two because it has that rank
on $U$.  At every parameter its fibre injects
into $H^0(X_s,\mathcal L_s)$ by the multiplication
exact sequence used in E.2.5.  Choose a frame
of $\mathcal V$ on the disc.  The common zero
locus of its two evaluated sections is closed
analytic in $\mathcal X$.  Its image is closed
analytic by properness and Remmert's theorem
\cite[Chapter~II, Theorem~8.8]{Demailly2012E}.
It misses $U$, so is a discrete analytic subset.
Remove it temporarily.

On the remaining connected parameter curve $S$
the sections define a proper holomorphic map
\[
 g:\mathcal X_S\longrightarrow\mathbb P^1\times S,
 \qquad \mathcal L=g^*\mathcal O(1).
 \tag{E.2.22.1}
\]
Each $g_s$ is nonconstant because its two sections
are independent.  Its fibres have pure dimension
two; the smooth total space and smooth target
then give flatness.  The fibres are connected:
the finite Stein factor is generically of degree
one, as it is so above $U$, and a finite
birational map to the normal target is an
isomorphism.  Its source is irreducible, so
there is no additional vertical component
in the Stein factor.  We retain
\[
 K_{\mathcal X_S/S}^{A}=g^*\mathcal O(D).
 \tag{E.2.22.2}
\]

\emph{Critical values and ordinary tori.}
The relative critical locus of $g$ is analytic
and proper over the target.  Its image has
finite intersection with each target fibre:
the nonconstant proper map $g_s$ to a curve
has finitely many critical values.  Thus the
reduced critical image is finite over $S$.
Its one-dimensional part has degree three,
as can be read over $U$; every such component
dominates $S$.  Remove the images of its
zero-dimensional components and its branch
and singular points.  These are further
proper analytic discrete subsets, absent
over the known open set.

A finite etale parameter cover now labels
the three critical sections.  Pointwise
classification after this cover descends;
images of later discrete exceptional sets
remain discrete because this cover is finite.
One may work on each connected component,
each of which meets the inverse image of $U$.
Send the labelled sections to $0,1,\infty$.
There are exactly these three critical values
on the current parameter domain.

The ordinary family over
$B^\circ\times S$, where
$B^\circ=\mathbb P^1\setminus\{0,1,\infty\}$,
is smooth and proper.  Its compact surface
fibres have the topology of the known two-tori
by Ehresmann's theorem and connectedness.
Adjunction gives integral $c_1(K_F)=0$,
since $F$ is a smooth fibre in a complex
six-sphere and its normal line is trivial.
These surfaces are complex tori.  To spell
out the classification step, $b_1(F)=4$
is even, so $F$ is K\"ahler.  The zero
canonical class excludes exceptional curves
and ruled surfaces.  On a K\"ahler surface,
a nonzero effective divisor cannot have
zero real Chern class; hence a nonzero
pluricanonical section has no zeros and
cannot give positive Kodaira dimension.
The compact-surface classification leaves
the Kodaira-dimension-zero list, in which
$b_1=4$ identifies a two-torus
\cite[Chapter~VI, Theorem~1.1 and Table~10]{BHPV2004E}.

\emph{The finite root charts.}
For any labelled critical section $\Gamma_j$,
every component of the Cartier divisor $g^*\Gamma_j$
dominates $S$.  A vertical divisor component would
be an entire fibre $X_s$, making $g_s$ constant.
Every dominating component meets $U$, where the
reduced fibre is irreducible.  Thus there is
exactly one such component, and its multiplicity
is the known one.  In particular
\[
 g^*\Gamma_1=(3k)\mathcal D_1,
 \qquad g^*\Gamma_2=(4l)\mathcal D_2,
 \qquad g^*\Gamma_0=\mathcal W,
 \tag{E.2.22.3}
\]
with reduced total hypersurfaces on the right.
For clarity, two distinct dominating components
could not both restrict to the known irreducible
reduced fibre.  Each is a Cartier prime in
the smooth total space, is torsion-free and
hence flat over the smooth parameter curve,
and has pure two-dimensional fibres.  Thus
its fibre over any parameter in $U$ contains
the whole known irreducible reduced surface.
At a general smooth point of
that fibre the known local equation has only
one reduced total branch.  Properness forces
every dominating component to meet that open
part of the known fibre, a contradiction.

For a finite multiplicity $m$, write locally
$q=a^m u$, with $a$ a reduced local equation
of $\mathcal D_j$ and $u$ a unit.  The normalization
of the cyclic root base change $r^m=q$ is
locally the disjoint union of the graphs
\[
 r=\zeta a u^{1/m},\qquad \zeta^m=1.
\]
It is therefore smooth and etale over the
actual total space, with free cyclic deck
action.  On $r\ne0$ its map to the root
disc times $S$ is already smooth.  Its
critical locus is supported over $r=0$,
is proper over $S$, and is empty over $U$.
Delete its discrete analytic image for
each of the two roots.  The resulting
proper smooth root families have connected
fibres by their Stein factors and are
families of two-tori by the same topology,
adjunction and surface argument above.

\emph{A finite analytic test for relative normal crossings.}
The total cusp $\mathcal W$ in (E.2.22.3)
is a reduced embedded hypersurface in
the smooth fourfold.  Take its finite
analytic normalization
\[
 \nu:\widetilde{\mathcal W}\longrightarrow
       \mathcal W\longrightarrow\mathcal X_S.
\]
Existence and finiteness are the analytic
normalization theorem
\cite[Chapter~II, Definition~7.9 and Theorem~7.12]{Demailly2012E}.
The normalization is proper over $S$.
First remove the proper images of its
singular locus, and then of the critical
locus of its map to $S$.  Each is analytic
and misses $U$.  On the resulting domain
the normalization is smooth over $S$.
Next remove the proper image of the
rank-drop locus of
\[
 T_{\widetilde{\mathcal W}/S}
       \longrightarrow\nu^*T_{\mathcal X_S/S}.
\]
This locus too is analytic and absent over
$U$.  The finite map $\nu$ is now an immersion.

For $r=2,3,4$, let $R_r$ be the locus
of ordered, pairwise distinct preimages
in the $r$-fold fibre product of
$\widetilde{\mathcal W}$ over $\mathcal X_S$.
Its partial diagonals are both open and
closed: a finite immersion is locally
an embedding, so its diagonal is open,
and separatedness makes it closed.
Consequently $R_r$ is closed analytic
in the proper fibre product and is
proper over $S$.

At a point of $R_r$ the $r$ local sheets
have relative conormal lines in the same
rank-three bundle $T^*_{\mathcal X_S/S}$.
On $R_2$ and $R_3$, delete the proper
images of the rank-drop loci of these
two or three lines.  Delete also the
image of $R_4$.  These are closed
analytic images and miss $U$, hence
are discrete.  There are only finitely
many such deletions.

At a remaining point the total cusp
has at most three local sheets, each
embedded smoothly and submersive over
$S$, with independent vertical normal
covectors.  Since the \emph{total}
$\mathcal W$ is a reduced embedded
hypersurface, its local ideal is exactly
the product of the distinct local
equations of those sheets.  No finite
subring or seminormality assumption
is being used.  Independence of the
vertical normals makes these equations,
the parameter on $S$, and the remaining
coordinates a holomorphic coordinate
system.  Thus
\[
 q=x,\qquad q=xy,\qquad\hbox{or}\qquad q=xyz
 \tag{E.2.22.4}
\]
after absorbing a unit.  These are
actual relative normal-crossing charts;
in particular the remaining cusp fibres
are reduced.  This proof did not assume
that normalization commutes with base
change before obtaining the charts.

\emph{The normalized toric pair and its action.}
The smooth proper normalization and
its conductor strata have the same
topology, incidence and normal degrees
as over $U$.  A further finite etale
cover can label these finitely many
strata.  The argument of E.2.18 applies
to every fibre: the six conductor
curves are a hexagon of $(-1)$-curves,
three alternate contractions give a
triangle in $\mathbb P^2$, and the
normalized pair is $(\mathrm{dP}_6,D)$.
Its relative logarithmic tangent bundle
is the evaluation pullback of a rank-two
bundle on the parameter curve.

In this evaluation description the
conductor descent maps are holomorphic
endomorphisms of that rank-two bundle.
They are constant along the compact
conductor curves.  They equal the identity
over $U$, where both toric logarithmic
fields descend.  The holomorphic identity
theorem on the connected parameter curve
makes them the identity everywhere.
This is the descent proof of E.2.19;
it requires a nonempty known open set,
not specifically a punctured disc.
Grauert evaluation, commutation on the
ordinary tori, completeness and the
implicit extension of the two primitive
toric periods then give the actual
marked cusp torus action locally on
the current parameter curve.  The
integral monodromy and triple-point
characters are constant by transport
and are the prescribed ones.

\emph{Continuation of the full period kernel.}
Fix an ordinary target point $t_*$.
The integral local system over
$B^\circ\times S$ has constant marked
peripheral representation, up to a
choice of marking: its matrices are
transported from those over $U$.
Parameter monodromy commutes with the
base representation.  It therefore
preserves the cusp invariant lattice
\[
 \ker(A_0-I)=\Z\widehat w\oplus\Z\widehat\delta.
\]
The wedge of the two corresponding
holomorphic period vectors at $t_*$
has a well-defined zero locus on $S$:
changing its integral basis multiplies
the wedge by a nonzero determinant.
It is nonzero over $U$, so that zero
locus is discrete analytic.  Remove
it as one final exceptional layer.

On a connected universal cover of
the remaining parameter curve, normalize
the Lie frame at $t_*$ by these two
columns.  The $\widehat\gamma$ column
then defines a holomorphic function
$c(s)$ by subtracting the fixed
reference value of $\beta^{\rm part}(t_*)$.
All other column identities at $t_*$
hold over a lift of $U$ and hence
everywhere by the identity theorem.

Now work on the universal cover of
the whole ordinary base times this
parameter cover.  The actual holomorphic
period kernel is a rank-two subbundle
of the complexified marked rank-four
lattice.  Compare it with
\[
 \ker\Pi_{c(s)}(t).
 \tag{E.2.22.5}
\]
The latter matrix always has complex
rank two, since its $\widehat w$ and
$\widehat\delta$ columns are the
standard basis.  These two holomorphic
Grassmannian maps agree over the
product with a lift of $U$.  They
therefore agree on the entire
connected ordinary cover.
This comparison does not assume
that the two actual columns were
independent at every $t$; equality
of the kernels proves it afterwards.
In a local actual Lie frame the
identity says that the full period
matrix is $W(t,s)\Pi_{c(s)}(t)$,
where $W$ is its two marked columns.
If $W$ had rank at most one, every
period column would have that rank,
contradicting the real rank four
of the compact torus lattice.

The actual ordinary and finite-root
fibres are compact tori.  The exact
period determinant is nonzero there.
Its sign is the negative sign on
$U$ and stays negative by connectedness
and nonvanishing.  Equality of the
period kernels also extends to the
smooth finite root charts; their
nonzero determinant preserves this
sign at the root point where the
maximum in R3 is attained.  Thus
$\operatorname{Im}c(s)+M<0$ by
the precise threshold of R3.
This places $c(s)$ in the
admissible interior domain, including
strictness at its boundary.  The
cusp multipliers are the corresponding
marked reference multipliers.  The
proof of E.2.20 now reconstructs
the actual cusp at each parameter:
relative normal-crossing transport
gives its known cover and incidence,
the action gives the same characters,
and adjacent-origin comparison gives
the monomial times unit extensions.
The proof works along any path from
$U$ and locally in the current
parameter; it does not need an
already known punctured neighbourhood
at the point being classified.

The finite affine lattice classes
and integral collar classes are
locally constant on this connected
domain, hence remain those over
$U$.  The additive comparison in
E.1 and E.2.21 identifies every
remaining actual threefold with
the reference member at $c(s)$.
The pointwise conclusion descends
from every finite parameter cover
used to label the critical or
conductor strata.

\emph{Filling the exceptional layers.}
We have made only finitely many
successive deletions.  At each stage
the deleted subset was closed
analytic and discrete on the
parameter domain then present.
Finite covers do not change that
property after taking their proper
images.  On the final domain all
fibres are known members of the
fixed pair.  At a point of the
last deleted set, choose a small
disc containing no other point of
that set and avoiding all earlier
deleted sets.  Its punctured disc
consists of known members.  E.2.21
identifies the central fibre as
well.  Fill the entire last layer
in this way, then the preceding
one, and continue in reverse order.
This is valid even when a later
layer accumulates at an earlier
deleted point.  Finally every
fibre over the original disc is
known.  No analytic curve selection
for an arbitrary open locus has
been assumed.
\end{proof}

% Owner: classification_e. Root and boundary_d full-source referee PASS.
\begin{theorem}[E.2.23 (Proved): separation under arbitrary deformation chains]
In a proper smooth holomorphic family of compact complex
manifolds over a connected reduced complex analytic base,
if one fibre is a sphere member of E.2.1 with coprime
multiplicities $(3k,4l)$, then every fibre is a constructed
member with that same pair.

Consequently two of the constructed sphere members are
deformation equivalent, allowing chains of proper smooth
families over connected complex analytic bases, if and
only if their ordered multiplicity pairs agree.  In
particular the distinct coprime pairs give distinct
deformation classes, and E.2.4 is resolved by this theorem.
\end{theorem}
\begin{proof}
Fix the pair and let $U$ be the locus of its constructed
members in the base $B$.  It is open by the actual local
miniversality of E.2.1 and E.2.2.  We prove that it is
closed without assuming a punctured analytic disc can
be selected inside an arbitrary open subset.

First consider any proper smooth family over a connected
Riemann surface.  If its known-pair locus is nonempty,
Ehresmann's theorem makes all its fibres diffeomorphic
to $S^6$, and that locus is open.  At any point of its closure choose a
coordinate disc meeting that locus.  The intersection
contains a nonempty open subset of the disc, so E.2.22
makes the entire disc known.  Thus the known-pair locus
is closed and equals the connected Riemann surface.
Only this disc propagation, not a special choice of
punctured disc, is used below.

Let $b\in\overline U\subset B$.  If $b\in U$ there
is nothing to prove.  Otherwise choose a small local
representative and an irreducible branch $T$ through
$b$ which meets $U$ arbitrarily near $b$.  There are
only finitely many such local branches.  Its dimension
is positive.  The local parametrization theorem gives
a finite proper projection
\[
 \rho:T\longrightarrow\Delta^d,
 \qquad \rho^{-1}(0)=\{b\},
 \tag{E.2.23.1}
\]
after shrinking the representative and translating
$b$ to the origin.  This is the finite projection
of \citet[Chapter~II, Theorem~4.19, pp.~95--96]{Demailly2012E}.
The isolated central fibre is part of choosing
the generic projection for the germ; the
Weierstrass representative is then taken small
enough to contain no other point over zero.

Choose $a\in U\cap T$ in this representative,
with $\rho(a)\ne0$.  Let $L$ be the complex
line through $0$ and $\rho(a)$, restricted
to the polydisc.  An irreducible component
$C$ of the reduced analytic set
$\rho^{-1}(L)$ which contains $a$ has
dimension one.  Indeed $d-1$ linear
equations give local dimension at least
one, while finiteness of $\rho$ gives
dimension at most one.  Its finite
proper image in $L$ is one-dimensional,
so is all of $L$.  Therefore $C$ contains
the unique point $b$ over zero.

Normalize $C$.  Its normalization
$\widetilde C$ is a connected Riemann
surface, finite over $C$; analytic
normalization exists and a normal
curve is smooth by
\cite[Chapter~II, Theorems~7.6 and~7.12]{Demailly2012E}.
It has points above both $a$ and $b$.
Pull back the given proper smooth
family to $\widetilde C$.  Its
known-pair locus contains the
inverse image of the open set
$U\cap C$ near $a$.  The preceding
Riemann-surface argument makes
every fibre on $\widetilde C$
known, in particular the fibre
above $b$.  Thus $b\in U$.

We have proved that $U$ is open
and closed.  Since $B$ is
connected and $U$ is nonempty,
$U=B$.  A nonreduced base
does not add fibre isomorphism
classes: restriction to its
reduction preserves both the
fibres and smoothness of the
family.  Thus the same conclusion
also applies to such a base
when it is admitted in the
definition of deformation
equivalence.

Applying this result to each
link in a finite deformation
chain preserves the pair.
Conversely, E.1 identifies
all admissible sphere data
with a fixed pair with the
connected period family,
using its holomorphic
comparisons and its period
lattice $2kl\Z$.
The admissible half-plane
and its punctured-disc
quotient are connected.
Thus all members with
that pair are deformation
equivalent.  This proves
both directions.
\end{proof}

% Owner: classification_e. Root and boundary_d full-source referee PASS.
\begin{corollary}[E.2.24 (Proved): separation for continuous complex structures]
Let $\mathcal I(S^6)$ be the space of smooth integrable
almost-complex structures on the fixed smooth sphere,
with its $C^\infty$ topology.  For a fixed coprime pair
$(3k,4l)$, let $\mathcal I_{k,l}$ consist of the structures
whose complex manifolds are biholomorphic to constructed
members with that pair.  Then $\mathcal I_{k,l}$ is open
and closed in $\mathcal I(S^6)$.

Consequently a continuous family of such structures
over any connected topological parameter space cannot
contain constructed members with two different pairs.
The same separation holds in the quotient spaces by
$\operatorname{Diff}_0(S^6)$ or by all orientation-preserving
diffeomorphisms, with their quotient topologies.
This statement does not identify the number of marked
connected components belonging to a single pair.
\end{corollary}
\begin{proof}
We use the classical Kuranishi local-completeness
theorem at an arbitrary $J_0\in\mathcal I(S^6)$,
including an initially unconstructed structure.
The complex orientation is locally constant, so the
oriented formulation applies on each such neighbourhood.
The precise $C^\infty$ version used here is
\citet[Section~3.2, Theorem~3.4 and Remark~3.6,
pp.~898--899]{Meersseman2019E}.  It supplies a connected
open neighbourhood $\mathcal U$ of $J_0$, a
finite-dimensional Kuranishi analytic space $K$
containing $J_0$, and a continuous retraction
$\Xi:\mathcal U\to K$ such that every $J\in\mathcal U$
is obtained from $\Xi(J)$ by a diffeomorphism isotopic
to the identity.  In that theorem the local product
map is first constructed in Sobolev completions;
elliptic regularity gives precisely this restriction
to smooth operators.  The topology on these operators
is the $C^\infty$ topology of
\cite[Section~2.1, pp.~883--884]{Meersseman2019E}.

The Kuranishi space carries its proper smooth
holomorphic deformation family.  We can take $K$
connected: in the stated local theorem it is the
continuous retract of the connected $\mathcal U$.
If necessary replace $K$ by its reduction, which
changes neither its points nor the complex manifolds
in its fibres.  Suppose $\mathcal U$ contains some
$J\in\mathcal I_{k,l}$.  Its Kuranishi representative
$\Xi(J)$ has the same complex manifold up to
biholomorphism.  E.2.23 applied to the connected
Kuranishi family therefore makes every point of
$K$ a member of the same pair.  Since all points
of $\mathcal U$ are biholomorphic to their
Kuranishi representatives, it follows that
$\mathcal U\subset\mathcal I_{k,l}$.

Thus at every $J_0$ the neighbourhood $\mathcal U$
is either contained in $\mathcal I_{k,l}$ or
disjoint from it.  This proves both openness and
closedness.  The inverse image of this set under
a continuous map from a connected parameter space
is then empty or the whole space, proving the
continuous-family assertion.  Finally
$\mathcal I_{k,l}$ is saturated under the indicated
diffeomorphism groups, because its definition uses
only the biholomorphism class.  Its image in either
quotient is open and closed by the definition of
the quotient topology.  No assertion that the
saturated set itself is connected is needed.
\end{proof}


